% ROZDZIAŁ 28 -- wzory miarowo-geometryczne, nie powtórzenie całek po podrozmaitościach.
\chapter{Wzór na pole powierzchni, krotności i koobszar}
\label{ch:area-coarea-28}

\section{Od wzoru parametrycznego do wzoru miarowego}
W skrypcie \emph{Geometria i topologia różniczkowa}
udowodniono wzór na całkę po kawałku gładkiej
podrozmaitości opisanej parametryzacją, która
pokrywa ten kawałek dokładnie raz. Tutaj rozważamy
\emph{odwzorowania nieiniektywne} oraz ich włókna.
Nowością są krotności obrazów, całkowanie po
przeciwobrazach punktów i formuły dla
odwzorowań Lipschitzowskich.

W tym rozdziale $\lambda_k$ oznacza miarę Lebesgue'a
na $\mathbb R^k$, a $\mathcal H^k$ -- miarę
Hausdorffa z Rozdziału~\ref{ch:hausdorff-26}.
Stosujemy konwencję z
Uwagi~\ref{rem:hausdorff-convention-26}: dla wymiaru
całkowitego $k\ge1$ wybieramy $c_k=\kappa_k^{-1}$,
tak aby $\mathcal H^k|_{\mathbb R^k}=\lambda_k$;
$\mathcal H^0$ jest miarą liczącą.
Dzięki temu poniższe wzory nie zakładają
nieudowodnionej nierówności izodiametrycznej.
Przy innej stałej normalizującej ich prawe strony
należy odpowiednio przeskalować.

\begin{definicja}[Jakobian wymiaru $k$]
\label{def:k-jacobian-28}
Niech $1\le k\le N$. Dla dowolnego odwzorowania
liniowego $A:\mathbb R^k\to\mathbb R^N$ definiujemy
\[
 J_k(A)=\sqrt{\det(A^{\mathsf T}A)}.
\]
Jeżeli $U\subseteq\mathbb R^k$ jest otwarty,
a $F:U\to\mathbb R^N$ jest różniczkowalne
w punkcie $x$, to kładziemy $J_kF(x)=J_k(DF(x))$.
W szczególności dotyczy to wszystkich punktów odwzorowania
$C^1$ oraz prawie wszystkich punktów odwzorowania
lokalnie Lipschitzowskiego. Macierz $A^{\mathsf T}A$ jest symetryczna i nieujemnie
określona, ponieważ $v^{\mathsf T}A^{\mathsf T}Av=|Av|^2$.
Ponadto $\ker(A^{\mathsf T}A)=\ker A$: jedna inkluzja
wynika z mnożenia przez $A^{\mathsf T}$, a druga z
poprzedniej równości. Zatem $J_k(A)=0$ wtedy i tylko
wtedy, gdy $\operatorname{rank}A<k$. Dla $k=N$
mamy $\det(A^{\mathsf T}A)=(\det A)^2$, więc
$J_k(A)=|\det A|$.
\end{definicja}

\begin{definicja}[Krotność odwzorowania]
\label{def:multiplicity-28}
Dla dowolnego odwzorowania $F:X\to Y$
i zbioru $E\subseteq X$ oraz dla każdego $y\in Y$ definiujemy
\[
 N(F,E,y)=\#\{x\in E:F(x)=y\}\in\mathbb N_0\cup\{+\infty\}.
\]
Symbol $\#$ oznacza liczbę elementów zbioru
(z wartością $+\infty$ dla zbiorów nieskończonych).
Pojawia się tu również wartość $+\infty$,
gdy włókno ma nieskończenie wiele punktów.
\end{definicja}

\section{Miara Hausdorffa gładkiego wykresu}
\begin{lemat}[Małe zniekształcenie wykresu]
\label{lem:graph-hausdorff-28}
Niech $1\le k\le N$, $V\subseteq\mathbb R^k$ będzie otwarty,
a $h:V\to\mathbb R^{N-k}$ funkcją klasy $C^1$.
Niech $\Gamma(z)=(z,h(z))$ oraz
$\nu(A)=\mathcal H^k(\Gamma(A))$
dla borelowskich $A\subseteq V$. Wtedy $\nu$
jest miarą Radona na $V$ i dla każdego
$z_0\in V$ zachodzi
\[
 \lim_{r\downarrow0}
 \frac{\nu(B(z_0,r))}{\lambda_k(B(z_0,r))}
 =\sqrt{\det\bigl(I_k+(Dh(z_0))^{\mathsf T}Dh(z_0)\bigr)}.
\]
\end{lemat}
\begin{proof}
Na każdej kuli domkniętej $K\subset V$ wykres $\Gamma$
jest Lipschitzowski, ponieważ $Dh$ jest ograniczona
na $K$. Zatem, na mocy
Twierdzenia~\ref{thm:hausdorff-lipschitz-26},
$\mathcal H^k(\Gamma(K))<\infty$, gdyż
$\mathcal H^k(K)=\lambda_k(K)$ jest skończona.
Z drugiej strony rzut ortogonalny $\pi$ na
pierwszych $k$ współrzędnych jest odwrotnością
$\Gamma$ na obrazie i jest $1$-Lipschitzowski.
Wykres $\Gamma(V)$ jest domknięty względem
$V\times\mathbb R^{N-k}$, więc jest borelowski w $\mathbb R^N$.
Dla borelowskiego $A\subseteq V$ mamy
$\Gamma(A)=\Gamma(V)\cap\pi^{-1}(A)$, zatem
$\Gamma(A)$ jest borelowski w $\mathbb R^N$.
Iniektywność $\Gamma$ zapewnia, że obrazy parami
rozłącznych zbiorów są parami rozłączne. Przeliczalna
addytywność $\mathcal H^k$ na zbiorach borelowskich
daje więc przeliczalną addytywność $\nu$.
Poprzednie oszacowanie zapewnia jej lokalną skończoność. Stosujemy twierdzenie o
regularności miar lokalnie skończonych na lokalnie
zwartych ośrodkowych przestrzeniach metrycznych
z Rozdziału~\ref{ch:radon-measures-20}.
Otrzymujemy, że $\nu$ jest miarą Radona,
bez zakładania z góry regularności miary
Hausdorffa na całym wykresie.

Ustalmy $z_0\in V$ i afiniczne odwzorowanie
$L(z)=(z,h(z_0)+Dh(z_0)(z-z_0))$.
Wybieramy $r>0$ tak, aby $\overline B(z_0,r)\subset V$,
i kładziemy
$\varepsilon_r=\sup_{w\in\overline B(z_0,r)}
\|Dh(w)-Dh(z_0)\|$. Z ciągłości $Dh$ wynika
$\varepsilon_r\to0$ dla $r\downarrow0$.
Dla $u,v\in B(z_0,r)$ odcinek między nimi
jest zawarty w tej kuli, więc podstawowe
twierdzenie rachunku całkowego daje
\[
 h(u)-h(v)-Dh(z_0)(u-v)
 =\int_{[0,1]}(Dh(v+t(u-v))-Dh(z_0))(u-v)
 \,\mathrm d\lambda_1(t).
\]
Norma prawej strony nie przekracza
$\varepsilon_r|u-v|$.
Ponieważ $|L(u)-L(v)|\ge|u-v|$,
odwzorowanie $\Gamma\circ L^{-1}$ między
odpowiednimi fragmentami wykresów jest
dwulipschitzowskie: dla $u,v$ z rozpatrywanej kuli
\[
 (1-\varepsilon_r)|L(u)-L(v)|
 \le|\Gamma(u)-\Gamma(v)|
 \le(1+\varepsilon_r)|L(u)-L(v)|.
\]
Dla dostatecznie małego $r$ mamy $\varepsilon_r<1$.
Z oszacowania miary Hausdorffa przez stałą
Lipschitza otrzymujemy
\[
 (1-\varepsilon_r)^k\mathcal H^k(L(B(z_0,r)))
 \le\nu(B(z_0,r))\le
 (1+\varepsilon_r)^k\mathcal H^k(L(B(z_0,r))).
\]
Obraz afinicznego odwzorowania $L$ jest
$k$-wymiarową płaszczyzną afiniczną. Wybierzmy izometrię
afiniczną tej płaszczyzny na $\mathbb R^k$ i oznaczmy przez
$B$ macierz części liniowej złożenia tej izometrii z $L$.
Wówczas $B^{\mathsf T}B=(DL)^{\mathsf T}DL=
I_k+Dh(z_0)^{\mathsf T}Dh(z_0)$.
Z niezmienniczości $\mathcal H^k$ względem izometrii,
kalibracji $\mathcal H^k=\lambda_k$ w $\mathbb R^k$
i liniowej zamiany zmiennych wynika
$\mathcal H^k(L(A))=|\det B|\lambda_k(A)$
dla borelowskiego $A\subseteq\mathbb R^k$.
Ponieważ $|\det B|=\sqrt{\det(B^{\mathsf T}B)}$,
jest to dokładnie $J_k(DL)\lambda_k(A)$.
Dzielimy przez $\lambda_k(B(z_0,r))$ i
przechodzimy do granicy.
\end{proof}

\begin{twierdzenie}[Miara na gładkim wykresie]
\label{thm:graph-area-28}
Niech $1\le k\le N$, $V\subseteq\mathbb R^k$ będzie otwarty,
$h:V\to\mathbb R^{N-k}$ będzie klasy $C^1$ i
$\Gamma(z)=(z,h(z))$. Dla każdej funkcji borelowskiej
$\psi:\Gamma(V)\to[0,+\infty]$ mamy
\[
 \int_{\Gamma(V)}\psi(y)\,\mathrm d\mathcal H^k(y)
 =\int_V\psi(\Gamma(z))
 \sqrt{\det(I_k+Dh(z)^{\mathsf T}Dh(z))}
 \,\mathrm d\lambda_k(z).
\]
To samo zachodzi dla funkcji całkowalnych.
\end{twierdzenie}
\begin{proof}
Z Lematu~\ref{lem:graph-hausdorff-28} wynika, że miara
$\nu(A)=\mathcal H^k(\Gamma(A))$ na zbiorach borelowskich
$A\subseteq V$ jest lokalnie skończona.
Pokrywamy $V$ przeliczalnie wieloma kulami o domknięciach
zawartych w $V$. Na każdej takiej kuli $\Gamma$ jest
Lipschitzowska. Jeśli $\lambda_k(A)=0$, to kalibracja
miary Hausdorffa w $\mathbb R^k$ i
Twierdzenie~\ref{thm:hausdorff-lipschitz-26} dają
zerową miarę obrazu przecięcia $A$ z każdą kulą.
Podaddytywność daje $\mathcal H^k(\Gamma(A))=0$,
czyli $\nu\ll\lambda_k$. Obie miary są
$\sigma$-skończone: każda z kul ma skończoną
miarę zarówno względem $\lambda_k$, jak i $\nu$.
Twierdzenie~\ref{thm:RN-sigfin-18} daje więc
gęstość $a=\mathrm d\nu/\mathrm d\lambda_k$.
Jest ona lokalnie całkowalna na $V$, gdyż
$\int_K a\,\mathrm d\lambda_k=\nu(K)<\infty$
dla każdego zwartego $K\subseteq V$.
Stosujemy Twierdzenie~\ref{thm:lebesgue-diff-25}
o różniczkowaniu całek lokalnie: na każdej kuli
o domknięciu zawartym w $V$ ograniczamy $a$ do
tej kuli i przedłużamy przez zero do funkcji
całkowalnej na $\mathbb R^k$. W każdym punkcie
wewnętrznym dostatecznie małe kule nie wychodzą
poza wybraną kulę. Przeliczalne pokrycie $V$
takimi kulami daje dla prawie każdego $z\in V$
\[
 a(z)=\lim_{r\downarrow0}
 \frac{\int_{B(z,r)}a\,\mathrm d\lambda_k}
 {\lambda_k(B(z,r))}
 =\lim_{r\downarrow0}\frac{\nu(B(z,r))}
 {\lambda_k(B(z,r))}=J_k\Gamma(z),
\]
gdzie ostatnia równość wynika z
Lematu~\ref{lem:graph-hausdorff-28}.
Tak otrzymujemy $\nu=J_k\Gamma\,\lambda_k$.
Dla borelowskiego $D\subseteq\Gamma(V)$,
z definicji $\nu$ mamy
$\nu(\Gamma^{-1}(D))=\mathcal H^k(D)$.
Stąd wzór na całki dla $\psi=\chi_D$, następnie
dla nieujemnych funkcji prostych przez liniowość,
a dla dowolnej nieujemnej borelowskiej $\psi$
przez rosnące przybliżenie funkcjami prostymi
i twierdzenie o zbieżności monotonicznej.
Dla funkcji całkowalnej rzeczywistej stosujemy
wzór do jej części dodatniej i ujemnej;
dla zespolonej do części rzeczywistej i urojonej.
\end{proof}

\section{Wzór na pole powierzchni dla immersji}
\begin{twierdzenie}[Wzór z krotnością dla immersji $C^1$]
\label{thm:immersion-area-28}
Niech $1\le k\le N$, $U\subseteq\mathbb R^k$
będzie otwarty, a $F:U\to\mathbb R^N$
immersją klasy $C^1$, tj.
$\operatorname{rank}DF(x)=k$ dla każdego $x\in U$.
Niech $E\subseteq U$ będzie borelowski,
a $f:E\to[0,+\infty]$ borelowska.
Wówczas funkcja
$y\mapsto\sum_{x\in E\cap F^{-1}(\{y\})}f(x)$
jest $\mathcal H^k$-mierzalna i
\begin{equation}\label{eq:area-immersion-28}
 \int_E f(x)J_kF(x)\,\mathrm d\lambda_k(x)
 =\int_{\mathbb R^N}\left(
 \sum_{x\in E\cap F^{-1}(\{y\})}f(x)\right)
 \,\mathrm d\mathcal H^k(y).
\end{equation}
Przy $f\equiv1$ na $E$ (czyli $f=\chi_E$ na $U$)
prawa strona redukuje się do
\[
 \int_{\mathbb R^N} N(F,E,y)\,\mathrm d\mathcal H^k(y).
\]
\end{twierdzenie}
\begin{proof}
Dla każdego $x\in U$ wybieramy $k$ składowych $F$,
których macierz pochodnej ma niezerowy wyznacznik w $x$.
Po permutacji współrzędnych celu oznaczmy rzut na te
składowe przez $\pi$. Twierdzenie o odwzorowaniu odwrotnym
daje otoczenie $U_x\subseteq U$, na którym
$\pi\circ F$ jest dyfeomorfizmem klasy $C^1$ na otwarty
zbiór w $\mathbb R^k$. W tych współrzędnych obraz $F(U_x)$
jest wykresem funkcji klasy $C^1$, a $F|_{U_x}$ jest osadzeniem.
Przeliczalność bazy pozwala wybrać przeliczalne
pokrycie takimi otoczeniami $U_j$. Tworzymy
rozłączny podział borelowski
$B_1=U_1$, $B_j=U_j\setminus\bigcup_{i<j}U_i$.

Na każdym $U_j$ korzystamy z opisanego wyżej wykresu
i dyfeomorfizmu $\phi_j=(\pi\circ F|_{U_j})^{-1}$
(z rzutem $\pi$ wybranym dla $U_j$).
Permutacja współrzędnych celu nie zmienia miary Hausdorffa
ani jakobianu $J_kF$.
Twierdzenie~\ref{thm:graph-area-28}, zastosowane
w tych współrzędnych, oraz klasyczny wzór
podstawienia na dziedzinie $\mathbb R^k$
dają dla borelowskich $A\subseteq B_j\cap E$
\[
 \int_A f(x)J_kF(x)\,\mathrm d\lambda_k(x)
 =\int_{F(A)}f((F|_{U_j})^{-1}(y))
 \,\mathrm d\mathcal H^k(y).
\]
Różniczka zmiany parametrów wprowadza dokładnie
ten czynnik, który przekształca macierz Grama
$F$ w macierz Grama wykresu; jest to tożsamość
$J_k(F\circ\phi)=J_kF(\phi)|\det D\phi|$
dla lokalnego dyfeomorfizmu $\phi$, wynikająca
z reguły łańcuchowej i równości
$\det((D\phi)^{\mathsf T}(DF)^{\mathsf T}DF D\phi)
=(\det D\phi)^2\det((DF)^{\mathsf T}DF)$,
gdzie $DF$ jest obliczane w punkcie $\phi(z)$.

Każdy $F(B_j\cap E)$ jest zbiorem borelowskim
w lokalnie osadzonym wykresie, a zatem również
w $\mathbb R^N$. Funkcja
$f\circ(F|_{U_j})^{-1}$ na $F(B_j\cap E)$
jest borelowska, gdyż odwrotność osadzenia jest ciągła.
Przedłużamy ją przez zero poza $F(B_j\cap E)$;
tak otrzymana funkcja jest borelowska na $\mathbb R^N$.
Sumujemy równości po $j$; po lewej używamy
rozłączności $B_j$, a po prawej twierdzenia
o zbieżności monotonicznej dla nieujemnych sum
częściowych. Nie wymaga ono $\sigma$-skończoności
miary Hausdorffa w przestrzeni otaczającej. Dla każdego $y$ suma
wkładów z kolejnych $B_j$ jest dokładnie sumą
po całym włóknie $E\cap F^{-1}(\{y\})$.
\end{proof}

\begin{przyklad}[Dlaczego trzeba liczyć krotności?]
Niech $F:(0,4\pi)\to\mathbb R^2$,
$F(t)=(\cos t,\sin t)$. Wówczas $J_1F=1$,
więc lewa strona \eqref{eq:area-immersion-28}
dla $f\equiv1$ wynosi $4\pi$.
Dla $y\in S^1\setminus\{(1,0)\}$ istnieje dokładnie
jeden $s\in(0,2\pi)$ z $F(s)=y$, a wszystkimi
przeciwobrazami w $(0,4\pi)$ są $s$ i $s+2\pi$.
Punkt $(1,0)$ ma jeden przeciwobraz $2\pi$.
Poza $S^1$ krotność wynosi zero.
Aby niezależnie obliczyć miarę okręgu, stosujemy
wzór immersji do różnowartościowego ograniczenia
$F|_{(0,2\pi)}$ z funkcją $1$. Otrzymujemy
$\mathcal H^1(S^1\setminus\{(1,0)\})=2\pi$.
Singleton ma zerową $\mathcal H^1$-miarę, więc
$\mathcal H^1(S^1)=2\pi$. Prawa strona dla
$(0,4\pi)$ wynosi zatem $2\cdot2\pi=4\pi$.
Bez krotności otrzymalibyśmy wynik dwa razy za mały.
\end{przyklad}

\section{Jakobian normalny i całkowanie po włóknach}
\begin{definicja}[Jakobian normalny]
\label{def:normal-jacobian-28}
Niech $n\ge m\ge1$. Dla dowolnego odwzorowania
liniowego $A:\mathbb R^n\to\mathbb R^m$ definiujemy
\[
 \operatorname{Jac}_m(A)=\sqrt{\det(AA^{\mathsf T})}.
\]
Dla otwartego $U\subseteq\mathbb R^n$ i odwzorowania
$F:U\to\mathbb R^m$ różniczkowalnego
w punkcie $x$ przyjmujemy
$\operatorname{Jac}_mF(x)=\operatorname{Jac}_m(DF(x))$.
Definicja obejmuje zatem odwzorowania $C^1$ i punkty
różniczkowalności odwzorowań Lipschitzowskich.
Ponieważ $w^{\mathsf T}AA^{\mathsf T}w=|A^{\mathsf T}w|^2$,
macierz $AA^{\mathsf T}$ jest nieujemnie określona,
a jej jądro jest równe $\ker A^{\mathsf T}$.
Jej wyznacznik jest dodatni dokładnie wtedy, gdy
$A$ ma rangę $m$. Stąd $\operatorname{Jac}_mF(x)=0$
dla rangi mniejszej niż $m$, a dla pełnej rangi
jest dodatni. Dla $m=1$ macierz $DF(x)$ jest
wierszem pochodnych cząstkowych, więc
$DF(x)DF(x)^{\mathsf T}=|\nabla F(x)|^2$.
Dla $n=m$ wyznacznik $AA^{\mathsf T}$ wynosi
$(\det A)^2$, zatem $\operatorname{Jac}_n(A)=|\det A|$.
\end{definicja}

\begin{lemat}[Tożsamość liniowo-algebraiczna koobszaru]
\label{lem:linear-coarea-28}
Niech $n\ge m\ge1$, $A:\mathbb R^n\to\mathbb R^m$
będzie surjektywnym odwzorowaniem liniowym, a
$G:\mathbb R^m\times\mathbb R^{n-m}\to\mathbb R^n$
odwracalnym odwzorowaniem liniowym, dla którego
$AG(y,z)=y$. Wówczas
\[
 \sqrt{\det(AA^{\mathsf T})}\,|\det DG|
 =\sqrt{\det((D_zG)^{\mathsf T}D_zG)}.
\]
\end{lemat}
\begin{proof}
Wybieramy ortonormalne rozkłady
$\mathbb R^n=(\ker A)^\perp\oplus\ker A$
i zapisujemy macierze względem tych baz.
Ograniczenie $A_0=A|_{(\ker A)^\perp}$
jest iniekcją, ponieważ jego jądro jest przecięciem
$\ker A$ z $(\ker A)^\perp$, czyli $\{0\}$.
Obie przestrzenie mają wymiar $m$, więc $A_0$
jest izomorfizmem. W tych bazach $A=(A_0\ \ 0)$.
Z warunku $AG(y,z)=y$ wynika, że pierwszy
wiersz blokowy macierzy $G$ to $(A_0^{-1}\ \ 0)$.
Macierz $G$ ma zatem postać
$\begin{psmallmatrix} A_0^{-1}&0\\B&C\end{psmallmatrix}$,
gdzie $C$ jest odwracalną macierzą kwadratową
rozmiaru $n-m$. Stąd
$|\det DG|=|\det A_0|^{-1}|\det C|$.
Ponieważ $AA^{\mathsf T}=A_0A_0^{\mathsf T}$,
lewy pierwiastek jest równy $|\det A_0|$.
Kolumny $D_zG$ leżą w $\ker A$, a ich macierz
w wybranej ortonormalnej bazie jest $C$,
więc prawy pierwiastek wynosi $|\det C|$.
Dla $n=m$ przyjmujemy wyznacznik pustej
macierzy równy $1$ i dowód pozostaje poprawny.
\end{proof}

\begin{twierdzenie}[Wzór koobszaru dla submersji]
\label{thm:coarea-submersion-28}
Niech $n\ge m\ge1$, $U\subseteq\mathbb R^n$
będzie otwarty, a $F:U\to\mathbb R^m$
submersją klasy $C^1$, tj.
$\operatorname{rank}DF(x)=m$ dla każdego $x\in U$.
Dla każdej borelowskiej funkcji $g:U\to[0,+\infty]$
funkcja całki wewnętrznej w poniższym wzorze jest
borelowska i zachodzi
\begin{equation}\label{eq:coarea-smooth-28}
 \int_U g(x)\operatorname{Jac}_mF(x)\,\mathrm d\lambda_n(x)
 =\int_{\mathbb R^m}\left(
 \int_{F^{-1}(\{y\})}g(x)\,\mathrm d\mathcal H^{n-m}(x)
 \right)\mathrm d\lambda_m(y).
\end{equation}
Wewnętrzna całka dotyczy miary Hausdorffa
na włóknie, które, o ile nie jest puste, jest
$(n-m)$-wymiarową podrozmaitością klasy $C^1$ (a dla
$n=m$ jest zliczaniem punktów). Jeżeli jedna ze stron
jest skończona dla $|g|$, równość zachodzi także
dla rzeczywistych i zespolonych funkcji całkowalnych.
\end{twierdzenie}
\begin{proof}
Gdy $n=m$, różniczka $DF(x)$ jest odwracalna
w każdym punkcie $U$. Twierdzenie o odwzorowaniu
odwrotnym daje pokrycie $U$ otwartymi zbiorami,
na których $F$ jest dyfeomorfizmem klasy $C^1$
na otwarty obraz. Wybieramy przeliczalne podpokrycie
$(U_j)$ i rozłączny borelowski podział
$B_j=U_j\setminus\bigcup_{i<j}U_i$.
Ponieważ $\operatorname{Jac}_nF=|\det DF|$,
zamiana zmiennych na $U_j$ daje
\[
 \int_{B_j}g(x)|\det DF(x)|\,\mathrm d\lambda_n(x)
 =\int_{F(B_j)}g((F|_{U_j})^{-1}(y))
 \,\mathrm d\lambda_n(y).
\]
Zbiory $F(B_j)$ są borelowskie, ponieważ
$F|_{U_j}$ jest homeomorfizmem na otwarty obraz.
Funkcje po prawej, przedłużone przez zero poza
$F(B_j)$, są borelowskie. Ich suma jest dla każdego
$y$ sumą $g$ po $F^{-1}(\{y\})$, gdyż z każdego
$B_j$ pochodzi co najwyżej jeden punkt włókna.
Ponieważ $\mathcal H^0$ jest miarą liczącą,
zbieżność monotoniczna sum częściowych daje wzór
oraz borelowskość całki po włóknie.
W dalszym dowodzie przyjmujemy $n>m$.

\emph{Krok 1: lokalne współrzędne włóknowe.}
Ustalmy $x_0\in U$. Ponieważ $DF(x_0)$ ma rangę $m$,
po permutacji współrzędnych źródła możemy zapisać
$x=(a,z)\in\mathbb R^m\times\mathbb R^{n-m}$
tak, aby minor $D_aF(x_0)$ był odwracalny.
Odwzorowanie $x\mapsto(F(x),z)$ jest klasy $C^1$,
a jego różniczka w $x_0$ ma postać
$\begin{psmallmatrix}D_aF&D_zF\\0&I_{n-m}\end{psmallmatrix}$.
Jej wyznacznik jest niezerowy. Twierdzenie o odwzorowaniu
odwrotnym daje otwarte otoczenie $V$ punktu $x_0$
i dyfeomorfizm odwrotny $G:W\to V$ klasy $C^1$,
gdzie $W\subseteq\mathbb R^m\times\mathbb R^{n-m}$
jest otwarty. Z definicji tej odwrotności
$F(G(y,z))=y$ dla $(y,z)\in W$.
Najpierw rozważamy funkcję $g$ znikającą poza $V$.

\emph{Krok 2: zamiana zmiennych i Fubini.}
Przez klasyczną zamianę zmiennych dostajemy
\[
 \int_V g(x)\operatorname{Jac}_mF(x)\,\mathrm d\lambda_n(x)
 =\int_W g(G(y,z))\operatorname{Jac}_mF(G(y,z))
 |\det DG(y,z)|\,\mathrm d\lambda_n(y,z).
\]
Ponieważ $DF(G)D_{(y,z)}G=(I_m\ \ 0)$,
Lemat~\ref{lem:linear-coarea-28} zastosowany
punktowo do różniczek daje
$\operatorname{Jac}_mF(G)|\det DG|=
J_{n-m}(z\mapsto G(y,z))$.
Przedłużamy nieujemną borelowską funkcję
$(y,z)\mapsto g(G(y,z))J_{n-m}(z\mapsto G(y,z))$
przez zero poza $W$. Ponieważ obie euklidesowe miary
borelowskie są $\sigma$-skończone, twierdzenie Tonellego
zapewnia borelowskość całki po $z$ i daje
\[
 \int_{\mathbb R^m}\left(\int_{W_y}
 g(G(y,z))J_{n-m}(z\mapsto G(y,z))
 \,\mathrm d\lambda_{n-m}(z)\right)\mathrm d\lambda_m(y),
\]
gdzie $W_y=\{z:(y,z)\in W\}$.
Dla ustalonego $y$ zbiór $W_y$ jest otwarty, a
$z\mapsto G(y,z)$ jest osadzeniem klasy $C^1$
o obrazie $V\cap F^{-1}(\{y\})$. Jest iniekcją
jako ograniczenie dyfeomorfizmu $G$, a jego
różniczka ma rangę $n-m$ jako ograniczenie
odwracalnego $DG$ do ostatnich $n-m$ współrzędnych.
Twierdzenie~\ref{thm:immersion-area-28} zastosowane
na $W_y$ do funkcji $z\mapsto g(G(y,z))$ daje
\[
 \int_{W_y}g(G(y,z))J_{n-m}(z\mapsto G(y,z))
 \,\mathrm d\lambda_{n-m}(z)
 =\int_{V\cap F^{-1}(\{y\})}g\,\mathrm d\mathcal H^{n-m}.
\]
Każda krotność na obrazie wynosi $1$; jeśli $W_y$
jest pusty, obie strony są zerowe. Ponieważ $g$
znika poza $V$, otrzymujemy lokalnie prawą stronę
\eqref{eq:coarea-smooth-28}.

\emph{Krok 3: sklejenie.}
Wybieramy przeliczalne pokrycie $U$ otoczeniami
$V_j$ jak w kroku 1 i rozłączny podział borelowski
$B_1=V_1$, $B_j=V_j\setminus\bigcup_{i<j}V_i$.
Zastępujemy $g$ przez $g\chi_{B_j}$ i stosujemy
lokalną równość. Dla każdego $y$ zbiory
$B_j\cap F^{-1}(\{y\})$ tworzą rozłączny
borelowski podział włókna, więc całka po włóknie
jest sumą lokalnych całek. Każda z nich jest
borelowską funkcją $y$ na mocy kroku 2;
ich suma również jest borelowska. Zbieżność
monotoniczna po obu stronach daje wzór globalny.

Jeżeli całka dla $|g|$ jest skończona, wzór dla
$|g|$ pokazuje, że całka po włóknie z $|g|$
jest skończona dla prawie wszystkich $y$.
Dla rzeczywistej funkcji $g$ stosujemy wzór
do $g^+$ i $g^-$ i odejmujemy skończone całki.
Dla zespolonej $g$ stosujemy ten argument
do części rzeczywistej i urojonej; ich moduły
są ograniczone przez $|g|$.
\end{proof}

\begin{wniosek}[Różniczkowanie poziomic]
\label{cor:coarea-level-28}
Dla $U\subseteq\mathbb R^n$ otwartego i
funkcji $F\in C^1(U;\mathbb R)$
bez punktów krytycznych oraz dowolnej
nieujemnej funkcji borelowskiej $g$ zachodzi
\[
 \int_U g(x)|\nabla F(x)|\,\mathrm d\lambda_n(x)
 =\int_{\mathbb R}\left(\int_{F^{-1}(\{t\})}
 g(x)\,\mathrm d\mathcal H^{n-1}(x)\right)
 \,\mathrm d\lambda_1(t).
\]
\end{wniosek}

\section{Zastosowanie: rozkład miary obrazowej}
\begin{twierdzenie}[Gęstość obrazu przez submersję]
\label{thm:pushforward-coarea-28}
Niech $n\ge m\ge1$, $U\subseteq\mathbb R^n$
będzie otwarty, $F:U\to\mathbb R^m$ submersją
klasy $C^1$, a $w:U\to[0,+\infty]$ funkcją
borelowską.
Niech $\nu=w\lambda_n$ i $\eta=F_*\nu$.
Wówczas $\eta\ll\lambda_m$, a jej gęstość jest
równa
\[
 q(y)=\int_{F^{-1}(\{y\})}
 \frac{w(x)}{\operatorname{Jac}_mF(x)}
 \,\mathrm d\mathcal H^{n-m}(x)
\]
dla $\lambda_m$-prawie wszystkich $y$,
z możliwością przyjmowania wartości $+\infty$.
\end{twierdzenie}
\begin{proof}
Dla borelowskiego $A\subseteq\mathbb R^m$
stosujemy wzór koobszaru do funkcji
\[
 g(x)=\frac{\chi_{F^{-1}(A)}(x)w(x)}
 {\operatorname{Jac}_mF(x)}.
\]
Mianownik jest dodatni na całym $U$, więc
jest to funkcja borelowska nieujemna.
Wynik ma postać
$\eta(A)=\int_A q(y)\,\mathrm d\lambda_m(y)$.
To dowodzi bezwzględnej ciągłości i istnienia
nieujemnej borelowskiej gęstości $q$, również gdy
$\eta$ nie jest $\sigma$-skończona. Jeżeli $\eta$
jest $\sigma$-skończona, $q$ jest jej pochodną
Radona--Nikodyma względem $\lambda_m$ w sensie
Twierdzenia~\ref{thm:RN-sigfin-18}.
\end{proof}

\begin{przyklad}[Poziomice promienia w pierścieniu]
Niech $U=\{x\in\mathbb R^2:a<|x|<b\}$,
gdzie $0<a<b$, i $F(x)=|x|^2$.
Wtedy $|\nabla F(x)|=2|x|$. Coarea daje
\[
 \int_U 2|x|\,\mathrm d\lambda_2(x)
 =\int_{a^2}^{b^2}\mathcal H^1(F^{-1}(\{t\}))
 \,\mathrm d\lambda_1(t)
 =\int_{a^2}^{b^2}2\pi\sqrt t\,\mathrm d\lambda_1(t)
 =\frac{4\pi}{3}(b^3-a^3).
\]
Jest to ilustracja całkowania po \emph{włóknach},
a nie kolejny sposób obliczania powierzchni sfery.
\end{przyklad}


\section{Twierdzenie Rademachera: skąd bierze się różniczka prawie wszędzie?}
\label{sec:rademacher-28}

Wzory geometryczne wymagają jakobianu, czyli różniczki. W przypadku
odwzorowań klasy $C^1$ istnieje ona wszędzie, natomiast odwzorowanie
Lipschitzowskie może nie być różniczkowalne w pewnych punktach.
Twierdzenie Rademachera dowodzi, że taki zbiór wyjątkowy ma zerową
miarę Lebesgue'a. Samo istnienie różniczki prawie wszędzie nie dowodzi
jednak ogólnych wzorów area i coarea: do ich uzasadnienia potrzebne są
jeszcze twierdzenia o pokrywaniu i mierzeniu obrazów zbiorów krytycznych.

\begin{definicja}[Różniczkowalność Frécheta]
\label{def:frechet-rademacher-28}
Niech $n,m\ge1$, $U\subseteq\mathbb R^n$ będzie otwarty,
$f:U\to\mathbb R^m$ i $x\in U$. Mówimy, że $f$ jest
\emph{różniczkowalna w sensie Frécheta} w $x$, jeżeli istnieje
liniowe $A:\mathbb R^n\to\mathbb R^m$ takie, że
\[
 \lim_{\substack{h\to0,\ h\ne0\\x+h\in U}}
 \frac{|f(x+h)-f(x)-Ah|}{|h|}=0.
\]
Takie $A$ jest jednoznaczne: jeżeli $A,B$ spełniają
warunek graniczny, to dla każdego wektora jednostkowego $v$
podstawienie $h=tv$ i przejście do granicy daje
$|(A-B)v|=0$. Piszemy $A=Df(x)$.
Odwzorowanie $f$ nazywamy \emph{lokalnie Lipschitzowskim},
jeżeli dla każdego $x\in U$ istnieje otwarte otoczenie
$V_x\subseteq U$ i stała $L_x\ge0$ takie, że
$|f(y)-f(z)|\le L_x|y-z|$ dla wszystkich $y,z\in V_x$.
\end{definicja}

\begin{lemat}[Jednowymiarowa różniczkowalność]
\label{lem:rademacher-line-28}
Niech $a<b$, a $f:[a,b]\to\mathbb R$ będzie funkcją
Lipschitzowską ze stałą $L$. Wtedy istnieje
$g\in L^\infty((a,b),\lambda_1)$ takie, że
$|g|\le L$ prawie wszędzie oraz
\[
 f(t)-f(a)=\int_{(a,t)}g(s)\,\mathrm d\lambda_1(s)
 \quad (t\in[a,b]).
\]
W konsekwencji $f'=g$ prawie wszędzie.
\end{lemat}
\begin{proof}
Ponieważ $f$ jest Lipschitzowska, jest absolutnie ciągła.
Dla przejrzystości odzyskamy gęstość bez powoływania się na
odwrotne podstawowe twierdzenie rachunku całkowego.
Połóżmy $p(t)=f(t)+Lt$. Dla $s<t$ mamy
$p(t)-p(s)\ge -L(t-s)+L(t-s)=0$, więc $p$ jest
niemalejąca i ciągła. Jej przyrosty na półotwartych
przedziałach określają dodatnią miarę borelowską
$\mu_p((s,t])=p(t)-p(s)$; istnienie tej miary wynika
z konstrukcji Lebesgue'a--Stieltjesa przez rozszerzenie
Carathéodory'ego (ciągłość $p$ zapewnia ciągłość
premiary z góry na przedziałach). Dla każdego przedziału
$[s,t]$ zachodzi
$0\le\mu_p((s,t])\le2L(t-s)$.
Nierówność przenosi się z przedziałów półotwartych na
ich skończone rozłączne sumy, czyli na algebrę generującą
zbiory borelowskie $(a,b]$. Obie miary są skończone na
$(a,b]$, więc rodzina zbiorów borelowskich $E$, dla których
$\mu_p(E)\le2L\lambda_1(E)$, jest zamknięta na granice
rosnących i malejących ciągów zbiorów.
Twierdzenie o klasie monotonicznej daje
$\mu_p\le2L\lambda_1$ na całej $\sigma$-algebrze borelowskiej
$(a,b]$. Dominacja przez $2L\lambda_1$ zapewnia brak atomów,
w tym w $b$. Obie miary na $(a,b]$ są skończone,
a $\mu_p\ll\lambda_1$, więc Twierdzenie~\ref{thm:RN-finite-18}
daje nieujemną całkowalną gęstość $q$.
Gdyby $\lambda_1(\{q>2L\})>0$, to dla pewnego
$j\ge1$ zbiór $E_j=\{q\ge2L+1/j\}$ miałby dodatnią
miarę. Wówczas
$\mu_p(E_j)=\int_{E_j}q\,\mathrm d\lambda_1
\ge(2L+1/j)\lambda_1(E_j)>2L\lambda_1(E_j)$,
wbrew dominacji. Stąd $0\le q\le2L$ prawie wszędzie.
Dla $a<t\le b$ mamy
\[
 f(t)-f(a)=p(t)-p(a)-L(t-a)
 =\int_{(a,t)}(q(s)-L)\,\mathrm d\lambda_1(s),
\]
ponieważ punkt $t$ ma miarę zero. Dla $t=a$
obie strony wynoszą zero. Funkcja $g=q-L$ spełnia
$|g|\le L$ prawie wszędzie, więc jest całkowalna
na $(a,b)$. Twierdzenie~\ref{thm:FTC-lebesgue-25}
zastosowane do tej funkcji daje $f'=g$ prawie wszędzie.
\end{proof}

\begin{lemat}[Słabe i kierunkowe pochodne funkcji Lipschitzowskiej]
\label{lem:rademacher-direction-28}
Niech $n\ge1$, $U\subseteq\mathbb R^n$ będzie otwarty,
a $f:U\to\mathbb R$ będzie lokalnie Lipschitzowska.
Istnieją mierzalne funkcje $g_i\in L^\infty_{\mathrm{loc}}(U)$,
$i=1,\ldots,n$, takie, że dla $\varphi\in C_c^1(U)$
\[
 \int_U f(x)\partial_i\varphi(x)\,\mathrm d\lambda_n(x)
 =-\int_U g_i(x)\varphi(x)\,\mathrm d\lambda_n(x).
\]
Dla każdego ustalonego $v\in\mathbb R^n$ pochodna
kierunkowa $\partial_v f(x)=\lim_{t\to0}
(f(x+tv)-f(x))/t$ istnieje dla prawie wszystkich
$x\in U$ i jest tam równa $\sum_i v_i g_i(x)$.
\end{lemat}
\begin{proof}
Przypuśćmy najpierw, że $f$ jest Lipschitzowska
ze stałą $L$ na domknięciu otwartego prostopadłościanu
$Q\Subset U$. Dla ustalonego $i$ oraz $t\in\mathbb Q\setminus\{0\}$
niech $q_t(x)=(f(x+te_i)-f(x))/t$ tam, gdzie $x+te_i\in Q$.
Zbiór punktów istnienia pochodnej $\partial_i f$ w $Q$
jest borelowski. Istotnie, warunek Cauchy'ego dla
$q_t(x)$ przy wymiernych $t\to0$ można zapisać jako
\[
 \forall l\ge1\ \exists j\ge1\
 \forall t,s\in\mathbb Q,\quad
 0<|t|,|s|<1/j\ \Longrightarrow\ |q_t(x)-q_s(x)|\le1/l,
\]
gdzie porównujemy tylko przyrosty z $x+te_i,x+se_i\in Q$.
Każdy warunek dotyczący ustalonych $t,s$ jest borelowski,
a wszystkie kwantyfikatory przebiegają zbiory przeliczalne.
Ponieważ $Q$ jest otwarty, a iloraz różnicowy jest ciągły
w rzeczywistym parametrze $t\ne0$, warunek ten jest
równoważny istnieniu skończonej granicy dla rzeczywistych
$t\to0$. Z Lematu~\ref{lem:rademacher-line-28} każdy
przekrój zbioru braku pochodnej w kierunku osi $i$ ma
jednowymiarową miarę zero. Stosując Tonellego do funkcji
charakterystycznej tego zbioru borelowskiego, otrzymujemy
jego $\lambda_n$-miarę zero (dla $n=1$ wystarcza sam lemat).
Na zbiorze istnienia pochodnej jest ona granicą
ilorazów różnicowych dla $t=1/j$, zdefiniowanych przez
zero tam, gdzie $x+e_i/j\notin Q$; dla każdego $x\in Q$
ostatecznie te przyrosty należą do $Q$.
Przypisując pochodnej wartość zero w pozostałych punktach,
otrzymujemy borelowską funkcję o module nie większym
niż $L$. Niech $\varphi\in C_c^1(Q)$.
Na każdym przekroju $Q$ równoległym do osi $i$
piszemy zmienną tej osi jako $t\in(a,b)$.
Niech $u(t)$ i $\psi(t)$ będą odpowiednimi
ograniczeniami $f$ i $\varphi$. Lemat
\ref{lem:rademacher-line-28} daje
$u(t)=u(a)+\int_{(a,t)}u'(s)\,\mathrm d\lambda_1(s)$.
Funkcja $\psi$ znika w otoczeniach obu końców,
gdyż $\varphi$ ma zwarty nośnik w $Q$.
Stąd $\int_{(a,b)}u(a)\psi'(t)\,\mathrm d\lambda_1(t)=0$.
Dla pozostałego składnika stosujemy Fubiniego:
całkowana funkcja jest bezwzględnie całkowalna,
bo $|u'|\le L$ prawie wszędzie, a $\psi'$ jest
ograniczona i przedział jest skończony. Otrzymujemy
\begin{align*}
 \int_{(a,b)}u(t)\psi'(t)\,\mathrm d\lambda_1(t)
 &=\int_{(a,b)}u'(s)
    \left(\int_{(s,b)}\psi'(t)\,\mathrm d\lambda_1(t)\right)
    \mathrm d\lambda_1(s)\\
 &=-\int_{(a,b)}u'(s)\psi(s)\,\mathrm d\lambda_1(s).
\end{align*}
Całkując tę równość po pozostałych współrzędnych,
stosujemy Fubiniego jeszcze raz; $f$ i $\partial_i f$
są ograniczone na $Q$, a $\varphi$ i jej pochodna
mają zwarty nośnik. Dostajemy
$\int_Q f\partial_i\varphi\,\mathrm d\lambda_n
=-\int_Q(\partial_i f)\varphi\,\mathrm d\lambda_n$.
Dla $n=1$ wystarcza już równość jednowymiarowa.
Wybieramy przeliczalne pokrycie $U$ wnętrzami takich
prostopadłościanów. Na przecięciach lokalne pochodne
cząstkowe są równe prawie wszędzie, bo obie są zwykłą
pochodną cząstkową tam, gdzie istnieje. Definiujemy $g_i$
na rozłącznych borelowskich częściach pokrycia, wybierając
odpowiedni lokalny reprezentant. Przeliczalność pokrycia
zapewnia zgodność prawie wszędzie na każdym prostopadłościanie.
Funkcja $g_i$ jest lokalnie istotnie ograniczona, ponieważ
każdy zwarty podzbiór $U$ pokrywa skończenie wiele wnętrz
tych prostopadłościanów. Dla $\varphi\in C_c^1(U)$ wybieramy
gładki podział jedności podporządkowany skończonemu pokryciu
jej zwartego nośnika tymi wnętrzami. Zapisujemy $\varphi$
jako skończoną sumę funkcji testowych o nośnikach w tych
wnętrzach i sumujemy lokalne tożsamości całkowania przez
części. Daje to żądaną tożsamość na całym $U$.

Dla $v=0$ obie strony żądanej równości są zerowe.
Dla ustalonego $v\ne0$ wybieramy ortonormalną bazę
z ostatnim wektorem $e=v/|v|$. Stosujemy dotychczas
udowodnioną część lematu w tych współrzędnych.
Otrzymujemy istnienie $\partial_e f$ prawie wszędzie,
jej lokalne istotne ograniczenie i tożsamość całkowania
przez części. Ortogonalna zmiana współrzędnych
zachowuje miarę Lebesgue'a. Ponieważ
$\partial_v f=|v|\partial_e f$, dla każdej
$\varphi\in C_c^1(U)$ mamy
\begin{align*}
 \int_U(\partial_v f)\varphi\,\mathrm d\lambda_n
 &=-\int_U f\sum_{i=1}^n v_i\partial_i\varphi\,\mathrm d\lambda_n\\
 &=\int_U\Bigl(\sum_{i=1}^n v_i g_i\Bigr)\varphi\,\mathrm d\lambda_n.
\end{align*}
Niech $w=\partial_v f-\sum_i v_i g_i$, gdzie
pochodną kierunkową przedłużamy przez zero na zbiór
braku pochodnej. Jest to funkcja lokalnie całkowalna
i $\int_U w\varphi\,\mathrm d\lambda_n=0$
dla każdej $\varphi\in C_c^1(U)$.
Aby uzasadnić jednoznaczność, wybieramy nieujemną
$\zeta\in C_c^\infty(B(0,1))$ o całce $1$
i kładziemy $\zeta_r(z)=r^{-n}\zeta(z/r)$.
Dla $x\in U$ oraz $\overline B(x,r)\subset U$
funkcja $y\mapsto\zeta_r(x-y)$ należy do $C_c^1(U)$,
więc jej całka z $w$ wynosi zero.
Dla każdego punktu Lebesgue'a $x$ funkcji $w$ mamy
\[
 |w(x)|\le \|\zeta\|_\infty r^{-n}
 \int_{B(x,r)}|w(y)-w(x)|\,\mathrm d\lambda_n(y).
\]
Prawa strona dąży do zera na mocy
Twierdzenia~\ref{thm:lebesgue-diff-25},
zastosowanego lokalnie do $w$. Prawie każdy punkt
jest punktem Lebesgue'a, a zatem $w=0$
prawie wszędzie. To daje
$\partial_v f=\sum_i v_i g_i$ prawie wszędzie.
\end{proof}

\begin{twierdzenie}[Rademacher]
\label{thm:rademacher-28}
Niech $n,m\ge1$, $U\subseteq\mathbb R^n$ będzie
otwarty, a $f:U\to\mathbb R^m$ lokalnie Lipschitzowska.
Wtedy istnieje zbiór borelowski $N\subseteq U$ o
$\lambda_n(N)=0$ taki, że $f$ jest różniczkowalna
w sensie Frécheta w każdym punkcie $U\setminus N$.
Ponadto w każdym punkcie różniczkowalności
$\|Df(x)\|\le\operatorname{Lip}(f;V)$ dla
każdego otoczenia $V$ punktu $x$, na którym $f$
jest Lipschitzowska.
\end{twierdzenie}
\begin{proof}
Najpierw dowodzimy różniczkowalności prawie wszędzie
dla $m=1$. Po jej udowodnieniu stosujemy ten wynik
do każdej z $m$ składowych $f_\ell$ odwzorowania
wektorowego. Poza skończoną sumą zbiorów zerowych
wszystkie $Df_\ell(x)$ istnieją. Macierz $A$ o wierszach
$Df_\ell(x)$ spełnia wtedy
\[
 \frac{|f(x+h)-f(x)-Ah|^2}{|h|^2}
 =\sum_{\ell=1}^m
 \frac{|f_\ell(x+h)-f_\ell(x)-Df_\ell(x)h|^2}{|h|^2}
 \longrightarrow0,
\]
co dowodzi różniczkowalności wektorowej.
Niech $g_1,\ldots,g_n$ będą słabymi pochodnymi z
Lematu~\ref{lem:rademacher-direction-28}. Rozważmy
przeliczalny zbiór kierunków $D=\mathbb Q^n\setminus\{0\}$.
Dla każdego $v\in D$ na zbiorze pełnej miary
istnieje $\partial_v f(x)=\sum_i v_i g_i(x)$.
Przecięcie tych przeliczalnie wielu zbiorów nadal
ma pełną miarę. Ustalmy punkt $x$ tego przecięcia,
w którym wszystkie $g_i(x)$ są skończone, i liniowe
odwzorowanie $A_xh=\sum_i g_i(x)h_i$.

Pokażemy, że $f(x+h)=f(x)+A_xh+o(|h|)$.
Ustalmy $\varepsilon>0$. Z gęstości racjonalnych
kierunków na sferze można wybrać \emph{skończenie}
wiele wektorów jednostkowych
$q_1,\ldots,q_M\in\{v/|v|:v\in D\}$, tak aby
każdy $u\in S^{n-1}$ był w odległości mniejszej
od $\delta$ od któregoś $q_j$, gdzie
$\delta>0$ wybierzemy później. Dla $h=tu$, $t=|h|$,
wybieramy taki $q_j$. Jeżeli $L$ jest stałą
Lipschitza $f$ w otoczeniu $x$, to
\begin{align*}
 |f(x+h)-f(x)-A_xh|
 &\le |f(x+tu)-f(x+tq_j)|\\
 &\quad+|f(x+tq_j)-f(x)-tA_xq_j|
       +|tA_x(q_j-u)|\\
 &\le t(L+\|A_x\|)\delta+o_j(t).
\end{align*}
Ponieważ lista kierunków $q_j$ jest skończona,
po podzieleniu przez $t$ ilorazy $o_j(t)/t$
zbiegają do zera \emph{jednostajnie względem $j$}.
Wybieramy $\delta$ tak, aby
$(L+\|A_x\|)\delta<\varepsilon$, a potem
dostatecznie małe $t$. Wynika stąd
$\limsup_{h\to0}|f(x+h)-f(x)-A_xh|/|h|\le\varepsilon$.
Dowolność $\varepsilon$ daje różniczkowalność
Frécheta z $Df(x)=A_x$.

Pozostaje uzasadnić borelowskość zbioru wyjątkowego.
Dla $A\in\mathbb Q^{m\times n}$, $r>0$ i $x\in U$
połóżmy
\[
 R_{A,r}(x)=\sup_{\substack{h\in\mathbb Q^n,\;0<|h|<r\\
                                    x+h\in U}}
 \frac{|f(x+h)-f(x)-Ah|}{|h|},
\]
przy czym supremum pustego zbioru przyjmujemy równe
zeru. Dla ustalonych $A,r$ funkcja $R_{A,r}$ jest
borelowska: supremum liczymy po przeliczalnie wielu
borelowskich ilorazach na otwartych dziedzinach
$\{x\in U:x+h\in U\}$.
Punkt $x$ jest punktem różniczkowalności wtedy i tylko wtedy,
gdy dla każdego $q\ge1$ istnieją
$A\in\mathbb Q^{m\times n}$ i $j\ge1$ takie, że
$R_{A,1/j}(x)<1/q$. Istotnie, jeśli $Df(x)$ istnieje,
wybieramy macierz wymierną $A$ z
$\|A-Df(x)\|<1/(3q)$, a następnie $j$ tak duże,
aby reszta dla $Df(x)$, podzielona przez $|h|$,
była nie większa niż $1/(3q)$ dla $0<|h|<1/j$.
Wtedy $R_{A,1/j}(x)\le2/(3q)<1/q$.

W drugą stronę dla każdego $q$ wybieramy odpowiednią
macierz $A_q$ i promień $1/j_q$. Dla $p,q\ge1$
i dowolnego niezerowego $v\in\mathbb Q^n$ wybieramy
dodatnie wymierne $t$ tak małe, aby $h=tv$ należało
do obu dopuszczalnych zakresów przyrostów. Odejmując
reszty dla $A_q$ i $A_p$, otrzymujemy
\[
 |(A_q-A_p)h|\le (1/q+1/p)|h|.
\]
Po podzieleniu przez $t$ i użyciu gęstości $\mathbb Q^n$
nierówność obowiązuje dla wszystkich $v\in\mathbb R^n$,
więc $\|A_q-A_p\|\le1/q+1/p$. Z zupełności
skończenie wymiarowej przestrzeni macierzy otrzymujemy
$A_q\to A$. Dla ustalonego $q$ i dostatecznie małych
wymiernych $h\ne0$ mamy
\[
 \frac{|f(x+h)-f(x)-Ah|}{|h|}
 \le1/q+\|A_q-A\|.
\]
Przez ciągłość względem $h\ne0$ ta nierówność obowiązuje
również dla rzeczywistych $h$ w tym samym otwartym
zakresie. Dla każdej dodatniej tolerancji wybieramy
$q$ tak duże, aby prawa strona była od niej mniejsza,
co dowodzi różniczkowalności z różniczką $A$.
Zatem zbiór punktów różniczkowalności ma postać
\[
 \bigcap_{q\ge1}\ \bigcup_{A\in\mathbb Q^{m\times n}}
 \ \bigcup_{j\ge1}\{x\in U:R_{A,1/j}(x)<1/q\}.
\]
Jest borelowski jako przeliczalne przecięcie
przeliczalnych sum zbiorów borelowskich. Jego dopełnienie ma zerową
miarę na mocy poprzedniego argumentu.
Jeśli $f$ jest Lipschitzowska na otoczeniu $V$ punktu
różniczkowalności $x$ ze stałą $L$, to dla wektora
jednostkowego $u$ i dostatecznie małych $t>0$
mamy $|f(x+tu)-f(x)|/t\le L$.
Przechodząc $t\downarrow0$, dostajemy
$|Df(x)u|\le L$. Biorąc supremum po $|u|=1$,
a następnie infimum po dopuszczalnych stałych $L$,
otrzymujemy $\|Df(x)\|\le\operatorname{Lip}(f;V)$.
\end{proof}

\begin{uwaga}[Dlaczego pochodne cząstkowe nie wystarczają?]
Na przykład $f(x,y)=|x|$ ma pochodne cząstkowe
w punktach z $x\ne0$, ale nie ma różniczki na prostej
$x=0$; zbiór ten jest zerowy dla $\lambda_2$.
Ponadto istnienie wszystkich dwustronnych pochodnych
kierunkowych w \emph{jednym} punkcie nie zapewnia
różniczkowalności Frécheta. Dla funkcji określonej przez
$f(0,0)=0$ i $f(x,y)=x^3/(x^2+y^2)$ poza początkiem
mamy $\partial_v f(0,0)=v_1^3/(v_1^2+v_2^2)$ dla
$v=(v_1,v_2)\ne0$ (iloraz różnicowy ma tę wartość
dla każdego $t\ne0$), a $\partial_0 f(0,0)=0$.
Nie jest to funkcja liniowa kierunku: wartości dla
$(1,0)$ i $(0,1)$ wynoszą $1$ i $0$, lecz dla $(1,1)$
otrzymujemy $1/2$. Istotna w dowodzie
Rademachera jest zgodność pochodnych kierunkowych
z jedną \emph{słabą} różniczką, a następnie
argument skończonej siatki kierunków.
\end{uwaga}

\begin{przyklad}[Wszystkie pochodne cząstkowe w punkcie, ale brak różniczki]
Niech $f:\mathbb R^2\to\mathbb R$ będzie zadana przez
$f(0,0)=0$ oraz $f(x,y)=xy/\sqrt{x^2+y^2}$
dla $(x,y)\ne(0,0)$. Poza początkiem ma ona
ograniczony gradient: jest jednorodna stopnia $1$,
a jej pochodne pierwszego rzędu są ciągłe na
okręgu jednostkowym i jednorodne stopnia $0$.
W związku z tym $f$ jest globalnie Lipschitzowska
(dla odcinka przechodzącego przez $0$ dzielimy go
na dwa odcinki). Obie pochodne cząstkowe w $(0,0)$
istnieją i są równe $0$, ponieważ $f(x,0)=f(0,y)=0$.
Jednak dla $h=(t,t)$, $t>0$, mamy
$f(t,t)=t/\sqrt2$, a więc
$|f(h)-f(0)-0\cdot h|/|h|=1/2$.
Nie istnieje zatem różniczka Frécheta w $(0,0)$.
Pokazuje to, dlaczego sam argument z istnienia
pochodnych cząstkowych prawie wszędzie jest
niewystarczający dla twierdzenia Rademachera.
\end{przyklad}

\begin{przyklad}[Różniczkowalność prawie wszędzie a brak gładkości]
Niech $F:\mathbb R^2\to\mathbb R^2$,
$F(x,y)=(x,y+|x|)$. Jest globalnie Lipschitzowskie,
choć nie jest klasy $C^1$. Dla $x>0$ i $x<0$
ma odpowiednio macierze $\begin{psmallmatrix}1&0\\1&1\end{psmallmatrix}$
i $\begin{psmallmatrix}1&0\\-1&1\end{psmallmatrix}$,
a na prostej $x=0$ nie ma różniczki. Mimo to
$|\det DF|=1$ prawie wszędzie. Wzór zamiany
zmiennych jest poprawny, ale jego uzasadnienie
w klasie Lipschitzowskiej nie może ograniczać się
do zwykłego twierdzenia o dyfeomorfizmach $C^1$.
\end{przyklad}

\begin{uwaga}[Związek z miarami wektorowymi]
Twierdzenie Rademachera wykorzystuje skończony wymiar
\emph{przestrzeni wartości}. Dla dowolnej przestrzeni
Banacha $X$ nie musi być prawdą, że każda funkcja
Lipschitzowska $[0,1]\to X$ ma pochodną prawie wszędzie.
Dokładną granicę wyznacza \emph{własność Radona--Nikodyma},
która zostanie zdefiniowana i scharakteryzowana w
Rozdziale~\ref{ch:vector-measures-30}.
\end{uwaga}

\section{Aproksymacja typu Łuzina--Whitneya w klasie $C^1$}
Aby przejść od wersji gładkich do odwzorowań Lipschitzowskich,
nie wystarczy różniczkowalność prawie wszędzie. Potrzebujemy
możliwości zastąpienia odwzorowania na dużych zbiorach
\emph{dokładnie} funkcją klasy $C^1$, wraz z jej różniczką.
Poniżej dowodzimy odpowiedniego faktu, nie przyjmując
jako gotowego twierdzenia o przedłużaniu Whitneya.

\begin{lemat}[Przedłużenie zgodnych danych pierwszego rzędu]
\label{lem:whitney-c1-28}
Niech $K\subseteq\mathbb R^n$ będzie niepustym zbiorem zwartym,
$a:K\to\mathbb R^N$ i
$A:K\to\mathcal L(\mathbb R^n,\mathbb R^N)$ funkcjami ciągłymi.
Załóżmy, że
\begin{equation}\label{eq:whitney-condition-28}
 \lim_{r\downarrow0}\ \sup_{\substack{x,y\in K,\;0<|x-y|\le r}}
 \frac{|a(y)-a(x)-A(x)(y-x)|}{|y-x|}=0,
\end{equation}
przy czym supremum pustego zbioru jest równe zeru.
Istnieje funkcja $\widetilde a\in C^1(\mathbb R^n;\mathbb R^N)$
taka, że $\widetilde a|_K=a$ i
$D\widetilde a(x)=A(x)$ dla $x\in K$.
\end{lemat}
\begin{proof}
Wprowadzamy dwa moduły ciągłości: $\omega(t)$ jest
supremum z \eqref{eq:whitney-condition-28} z $r=t$,
a $\alpha(t)=\sup\{\|A(x)-A(y)\|:x,y\in K,
|x-y|\le t\}$. Obydwa dążą do zera dla $t\downarrow0$.

\emph{Krok 1: kostki otwartego dopełnienia.}
Zbiór $\mathbb R^n\setminus K$ rozkładamy na kostki
diadyczne $Q$ o rozłącznych wnętrzach, dobierając
kostki maksymalne (względem włączenia w rodzinie
kostek diadycznych) spośród tych, dla których
$20\sqrt n\,\ell(Q)\le\operatorname{dist}(Q,K)$.
Tutaj $\ell(Q)$ oznacza długość boku kostki.
Używamy domkniętych kostek standardowej siatki
diadycznej, tj. zbiorów $2^{-j}(b+[0,1]^n)$ dla
$j\in\mathbb Z$ i $b\in\mathbb Z^n$.
Dwie kostki tej siatki o przecinających się wnętrzach
są zagnieżdżone: w każdej współrzędnej końce przedziału
o większej długości należą do siatki wyznaczonej przez
mniejszą długość, więc wnętrze mniejszego przedziału
nie może przecinać końca większego. Stąd różne
maksymalne kostki mają rozłączne wnętrza.
Dla $x\notin K$ kładziemy $\delta_x=\operatorname{dist}(x,K)>0$.
Jeśli kostka $Q$ zawiera $x$ i
$21\sqrt n\,\ell(Q)\le\delta_x$, to
$\operatorname{dist}(Q,K)\ge\delta_x-\sqrt n\,\ell(Q)
\ge20\sqrt n\,\ell(Q)$.
Każdy punkt dopełnienia należy więc do dostatecznie
małej kostki spełniającej nierówność. Kolejne kostki
macierzyste zawierają ten sam punkt $x$, więc ich odległość
od $K$ nie przekracza $\operatorname{dist}(x,K)$, natomiast
ich boki rosną bez ograniczenia. Każda dopuszczalna kostka
ma więc maksymalnego przodka: jeżeli kostka nie spełnia
warunku, to jej więksi przodkowie też go nie spełniają,
ponieważ odległość od $K$ nie rośnie, a bok rośnie.
Maksymalne kostki pokrywają zatem całe dopełnienie $K$,
a każda z nich leży w tym dopełnieniu dzięki dodatniej
odległości od $K$. Kostka macierzysta $Q'$ kostki
maksymalnej $Q$ nie spełnia nierówności, skąd
\[
 \operatorname{dist}(Q,K)
 \le\operatorname{dist}(Q',K)+\operatorname{diam}Q'
 <42\sqrt n\,\ell(Q).
\]
Oznaczmy przez $Q^*$ współśrodkową kostkę o boku
$(9/8)\ell(Q)$. Dla $x\in Q^*$ mamy
\[
 19\sqrt n\,\ell(Q)\le\operatorname{dist}(x,K)
 \le44\sqrt n\,\ell(Q).
\]
Dolne oszacowanie wynika z odległości $x$ od $Q$,
nie większej niż $\sqrt n\,\ell(Q)/16$; górne
wynika z poprzedniej nierówności i średnicy $Q^*$.
Zatem gdy $Q^*\cap R^*\ne\varnothing$, stosunek boków
$Q,R$ zawiera się w $[19/44,44/19]$.
Dla ustalonego $x\notin K$ wszystkie kostki, dla których
$x\in Q^*$, mają boki porównywalne z
$\delta=\operatorname{dist}(x,K)>0$ i mieszczą się
w kuli $B(x,C_n\delta)$. Ich wnętrza są rozłączne,
a objętość każdej jest nie mniejsza niż $c_n\delta^n$.
Porównanie z objętością tej kuli ogranicza ich liczbę
przez stałą zależną tylko od $n$. Ten sam argument na
kuli $B(x,\delta/4)$, na której odległość od $K$ pozostaje
między $3\delta/4$ a $5\delta/4$, dowodzi lokalnej
skończoności rodziny powiększonych kostek w dopełnieniu $K$.

Dla każdej kostki wybieramy $p_Q\in K$ minimalizujący
odległość od jej środka; minimum istnieje dzięki
niepustości i zwartości $K$. Wybieramy $C^\infty$-funkcje
$\psi_Q\ge0$, równe $1$ na $Q$, o nośnikach zawartych
w kostkach współśrodkowych o boku $(9/8)\ell(Q)$,
i spełniające $\|D\psi_Q\|\le C_n/\ell(Q)$.
Dokładniej, wybieramy nieujemną funkcję
$\theta\in C_c^\infty((-9/16,9/16))$, równą $1$
na $[-1/2,1/2]$ i nie większą niż $1$.
Dla środka $c_Q$ kostki $Q$ kładziemy
$\psi_Q(x)=\prod_{i=1}^n\theta((x_i-(c_Q)_i)/\ell(Q))$.
Reguła łańcuchowa daje
$\|D\psi_Q\|\le\sqrt n\|\theta'\|_\infty/\ell(Q)$.
Suma $\sum_Q\psi_Q$ jest lokalnie skończona i nie
mniejsza niż $1$ na $\mathbb R^n\setminus K$.
Oznaczmy $S(x)=\sum_R\psi_R(x)$. Funkcja $S$
jest gładka na dopełnieniu $K$ i $S\ge1$.
Funkcje $\varphi_Q=\psi_Q/S$ są nieujemne, gładkie,
mają nośniki w $Q^*$ i sumują się do $1$.
Dla $x\in Q^*$ wszystkie kostki $R$, których składnik
może mieć niezerową pochodną w $x$, mają boki
porównywalne z $\ell(Q)$, a ich liczba jest ograniczona
przez stałą zależną tylko od $n$. Stąd
$\|DS(x)\|\le C_n/\ell(Q)$. Z reguły ilorazu,
$S\ge1$ oraz $0\le\psi_Q\le1$ wynika
\[
 \|D\varphi_Q(x)\|
 \le\frac{\|D\psi_Q(x)\|}{S(x)}
 +\frac{\psi_Q(x)\|DS(x)\|}{S(x)^2}
 \le C_n/\ell(Q).
\]
Poza $Q^*$ pochodna jest zerowa, więc oszacowanie
obowiązuje w całym dopełnieniu $K$.

\emph{Krok 2: sklejanie przybliżeń afinicznych.}
Dla $x\notin K$ definiujemy
\[
 P_Q(x)=a(p_Q)+A(p_Q)(x-p_Q),\qquad
 \widetilde a(x)=\sum_Q\varphi_Q(x)P_Q(x),
\]
a dla $x\in K$ przyjmujemy $\widetilde a(x)=a(x)$.
Niech $x\notin K$, $\delta=\operatorname{dist}(x,K)$,
a $p\in K$ będzie punktem najbliższym $x$.
Dla każdej kostki, dla której $x\in Q^*$,
mamy $\ell(Q)\asymp\delta$. Ponieważ $p_Q$ jest
punktem najbliższym środkowi $c_Q$, mamy
$|c_Q-p_Q|\le|c_Q-p|\le |c_Q-x|+\delta$.
Zatem $|x-p_Q|\le 2|x-c_Q|+\delta\le C_n\delta$
i $|p_Q-p|\le |p_Q-x|+\delta\le C_n\delta$.
Rozpisujemy różnicę przybliżeń afinicznych:
\begin{align*}
 P_Q(x)-a(p)-A(p)(x-p)
 &=a(p_Q)-a(p)-A(p)(p_Q-p)\\
 &\quad+(A(p_Q)-A(p))(x-p_Q).
\end{align*}
Pierwszy składnik ma normę nie większą niż
$|p_Q-p|\omega(|p_Q-p|)$ (jest zerowy dla $p_Q=p$),
a drugi nie większą niż
$\alpha(|p_Q-p|)|x-p_Q|$.
Z monotoniczności obu modułów otrzymujemy
\[
 |P_Q(x)-[a(p)+A(p)(x-p)]|
 \le C_n\delta\,[\omega(C_n\delta)+\alpha(C_n\delta)],
\]
zaś $\|A(p_Q)-A(p)\|\le\alpha(C_n\delta)$.
Ponieważ $\sum_Q\varphi_Q=1$, pierwsze z poniższych
oszacowań otrzymujemy przez sumowanie poprzedniego
oszacowania z wagami $\varphi_Q(x)$.
Ponadto $\sum_QD\varphi_Q=0$, więc w ustalonym punkcie $x$
(wybór $p$ jest tu ustalony, nie różniczkujemy go) mamy
\begin{align*}
 D\widetilde a(x)-A(p)
 &=\sum_Q\varphi_Q(x)(A(p_Q)-A(p))\\
 &\quad+\sum_Q [P_Q(x)-a(p)-A(p)(x-p)]D\varphi_Q(x).
\end{align*}
Druga suma oznacza sumę odwzorowań liniowych
$h\mapsto [P_Q(x)-a(p)-A(p)(x-p)]D\varphi_Q(x)h$.
Liczba niezerowych składników jest jednostajnie ograniczona,
a $\|D\varphi_Q(x)\|\le C_n/\delta$ dla uczestniczących
kostek, zatem otrzymujemy
\begin{align*}
 |\widetilde a(x)-a(p)-A(p)(x-p)|
 &\le C_n\delta[\omega(C_n\delta)+\alpha(C_n\delta)],\\
 \|D\widetilde a(x)-A(p)\|
 &\le C_n[\omega(C_n\delta)+\alpha(C_n\delta)].
\end{align*}
Stałe mogą również zależeć od wyboru funkcji
odcinających, ale nie od $x$.

Gdy $x\to p_0\in K$ od zewnątrz, mamy
$\delta\le|x-p_0|$ i $|p-p_0|\le2|x-p_0|$.
Stąd $p\to p_0$ i $\delta\to0$. Drugie oszacowanie
oraz ciągłość $A$ dają $D\widetilde a(x)\to A(p_0)$.
Dla pierwszego rozwinięcia piszemy
\begin{align*}
 \widetilde a(x)-a(p_0)-A(p_0)(x-p_0)
 &=\widetilde a(x)-a(p)-A(p)(x-p)\\
 &\quad+a(p)-a(p_0)-A(p_0)(p-p_0)\\
 &\quad+(A(p)-A(p_0))(x-p).
\end{align*}
Po podzieleniu przez $|x-p_0|$ normy kolejnych
składników są ograniczone odpowiednio przez
$C_n[\omega(C_n\delta)+\alpha(C_n\delta)]$,
$2\omega(2|x-p_0|)$ i $\alpha(2|x-p_0|)$.
Wszystkie trzy ograniczenia dążą do zera,
co daje żądany rozwój z resztą $o(|x-p_0|)$.
Ten sam warunek bezpośrednio daje taki rozwój,
gdy $x\in K$. Zatem $\widetilde a$ jest
różniczkowalna także na $K$, jej różniczka jest
tam równa $A$ i jest ciągła na całej przestrzeni.
\end{proof}

\begin{lemat}[Zwarte fragmenty klasy $C^1$]
\label{lem:lusin-whitney-28}
Niech $U\subseteq\mathbb R^n$ będzie otwarty,
$F:U\to\mathbb R^N$ lokalnie Lipschitzowski,
$Q\subseteq U$ zwarty i $\varepsilon>0$.
Istnieją zbiór zwarty $K\subseteq Q$ oraz
$F_K\in C^1(\mathbb R^n;\mathbb R^N)$ takie, że
\[
 \lambda_n(Q\setminus K)<\varepsilon,\qquad
 F_K(x)=F(x),\quad DF_K(x)=DF(x)\quad(x\in K).
\]
Różniczkę $DF$ rozumiemy tylko w punktach,
w których istnieje; wszystkie punkty $K$ są
takimi punktami. Ponadto istnieje przeliczalna rodzina zbiorów zwartych
$K_j\subseteq U$, z których każdy ma opisane wyżej
przedłużenie i dla której
$\lambda_n(U\setminus\bigcup_jK_j)=0$.
\end{lemat}
\begin{proof}
Niech $D\subseteq U$ będzie borelowskim zbiorem
punktów różniczkowalności $F$. Na mocy
Twierdzenia~\ref{thm:rademacher-28} dopełnienie
$D$ jest zerowe. Funkcja $x\mapsto DF(x)$ jest
borelowska na $D$, ponieważ jej współczynniki
są granicami ilorazów różnicowych w kierunkach
wektorów bazowych.

Dla $j,l\ge1$ oznaczmy przez $H_{j,l}$ zbiór
punktów $x\in D\cap Q$, dla których
\[
 |F(x+h)-F(x)-DF(x)h|\le j^{-1}|h|
\]
dla każdego $h\in\mathbb Q^n$ spełniającego
$0<|h|<l^{-1}$ i $x+h\in U$. Jest to zbiór
borelowski, bo warunek sprawdzamy dla
przeliczalnej liczby przyrostów. Z różniczkowalności wynika
$\bigcup_lH_{j,l}=D\cap Q$, przy czym dla ustalonego
$j$ zbiory $H_{j,l}$ rosną wraz z $l$.
Ponieważ $Q$ jest zwarty, $\lambda_n(Q)<\infty$,
więc ciągłość miary z dołu pozwala dla każdego $j$
wybrać $l_j$ tak duże,
aby $\lambda_n((D\cap Q)\setminus H_{j,l_j})
<\varepsilon2^{-j-2}$. Zbiór
$G=\bigcap_{j\ge1}H_{j,l_j}$ spełnia
\[
 \lambda_n(Q\setminus G)
 \le\lambda_n(Q\setminus D)
 +\sum_{j\ge1}\lambda_n((D\cap Q)\setminus H_{j,l_j})
 \le\varepsilon/4<\varepsilon/2.
\]
Z wewnętrznej regularności miary Lebesgue'a
wybieramy zwarty $K_0\subseteq G$ z
$\lambda_n(G\setminus K_0)<\varepsilon/4$.
Ograniczenie $DF$ do $K_0$ jest borelowskim
odwzorowaniem o wartościach w skończenie wymiarowej
przestrzeni macierzy, a miara Lebesgue'a ograniczona
do $K_0$ jest skończona i regularna. Stosujemy Twierdzenie~\ref{thm:lusin-10}
do tego odwzorowania (lub do wszystkich jego
$Nn$ współczynników, wybierając dla każdego zbiór zwarty
ze stratą mniejszą niż $\varepsilon/(4Nn)$ i biorąc
ich skończone przecięcie). Otrzymujemy zwarty
$K\subseteq K_0$, na którym $DF|_K$ jest ciągła
i $\lambda_n(K_0\setminus K)<\varepsilon/4$.
W konsekwencji $\lambda_n(G\setminus K)<\varepsilon/2$. Jeśli $K=\varnothing$,
przyjmujemy $F_K\equiv0$ i pierwsza teza jest
trywialna; w dalszym ciągu zakładamy $K\ne\varnothing$.

Warunek w definicji $H_{j,l_j}$, przez ciągłość
$F$ względem przyrostu $h$, pozostaje prawdziwy
również dla rzeczywistych $h$ o tej samej
ograniczonej długości. Stąd dla każdego $j$ i
$x,y\in K$ o $0<|x-y|<l_j^{-1}$ zachodzi
\[
 |F(y)-F(x)-DF(x)(y-x)|\le j^{-1}|x-y|.
\]
Dla każdej zadanej dodatniej tolerancji wybieramy $j$
tak, aby $1/j$ było od niej mniejsze; dla $r<1/l_j$
supremum ilorazów w warunku Whitneya nie przekracza
$1/j$. Ponadto $a=F|_K$ jest ciągłe dzięki lokalnej
Lipschitzowskości $F$, a $A=DF|_K$ jest ciągłe
z wyboru $K$. Zatem dane $a,A$ spełniają wszystkie
założenia Lematu~\ref{lem:whitney-c1-28}.
Jego zastosowanie daje żądane przedłużenie.

Aby uzyskać ostatnią tezę, wybieramy przeliczalne
pokrycie $U$ zwartymi kostkami $Q_s\Subset U$.
Dla każdego $s$ i $j$ wykonujemy konstrukcję
z $\varepsilon=1/j$ na $Q_s$, uzyskując
$K_{s,j}$. Ponieważ
$\lambda_n(Q_s\setminus\bigcup_jK_{s,j})\le1/j$
dla każdego $j$, miara ta jest zerowa.
Przeliczalna suma pozostałych zbiorów zerowych
ma miarę zero, co kończy dowód.
\end{proof}

\section{Dwa oszacowania dla zbiorów krytycznych}
\begin{lemat}[Obraz zbioru o niedostatecznej randze]
\label{lem:critical-area-null-28}
Niech $1\le k\le N$, $V\subseteq\mathbb R^k$ będzie
otwarty, $G\in C^1(V;\mathbb R^N)$, a
$K\subseteq V$ zwarty. Jeżeli
$\operatorname{rank}DG(x)<k$ dla wszystkich $x\in K$,
to $\mathcal H^k(G(K))=0$.
\end{lemat}
\begin{proof}
Jeśli $K=\varnothing$, teza wynika z $\mathcal H^k(\varnothing)=0$.
W przeciwnym razie wybieramy $\rho>0$ tak, aby
$K_\rho=\{z:\operatorname{dist}(z,K)\le\rho\}$ był
zwartym podzbiorem $V$. Istnieje takie $\rho$, ponieważ
$K$ jest zwarty i $V$ otwarty. Niech $M$ będzie
ograniczeniem normy $DG$ na $K_\rho$. Ustalmy
$0<\eta\le1$. Z jednostajnej ciągłości $DG$ na $K_\rho$
wybieramy $r>0$ dostatecznie małe, aby $\sqrt k\,r<\rho$
i aby dla punktów tego otoczenia odległych o co najwyżej
$\sqrt k\,r$ różnica ich różniczek miała normę co najwyżej
$\eta$. Warunki te zachodzą również dla wszystkich
mniejszych dodatnich $r$.

Rozpatrujemy domknięte kostki siatki o boku $r$,
przecinające $K$. Dla każdej takiej kostki $Q$
wybieramy $x_Q\in K\cap Q$. Dla $x\in K\cap Q$
cały odcinek $[x_Q,x]$ leży w $K_\rho$, więc
\begin{align*}
 G(x)-G(x_Q)-DG(x_Q)(x-x_Q)
 &=\int_0^1[DG(x_Q+t(x-x_Q))-DG(x_Q)]\\
 &\qquad\cdot(x-x_Q)\,\mathrm d\lambda_1(t).
\end{align*}
Norma tej różnicy nie przekracza $\eta\sqrt k\,r$.
Niech $P_Q=\operatorname{im}DG(x_Q)$ oraz
$j=\dim P_Q\le k-1$. Po odjęciu $G(x_Q)$
obraz $G(K\cap Q)$ leży w $\eta\sqrt k\,r$-otoczeniu
kuli w $P_Q$ o promieniu $M\sqrt k\,r$.

Dla $j\ge1$ wybieramy ortonormalne współrzędne w $P_Q$.
Kula ta mieści się w $j$-wymiarowej kostce o boku
$2M\sqrt k\,r$. Dzielimy ją siatką o boku $\eta r$.
Liczba przecinających ją kostek nie przekracza
$(2M\sqrt k/\eta+2)^j\le C_{k,M}\eta^{-j}$.
Każda z nich ma średnicę $\sqrt j\,\eta r$.
Jeśli $j=0$, część liniowa jest zerowa i wystarcza
jeden punkt $0\in P_Q$, czyli ten sam szacunek
liczby zbiorów przy wykładniku $j=0$.
Po powiększeniu każdej kostki do kuli w $\mathbb R^N$
o promieniu $(\sqrt k+\sqrt j)\eta r$ uzyskujemy
pokrycie przesuniętego obrazu. Ponieważ $j\le k-1$
i $\eta\le1$, koszt tego pokrycia nie przekracza
$C_{k,N,M}\eta^{-j}(\eta r)^k
\le C_{k,N,M}\eta r^k$.

Zwartość $K$ zapewnia, że leży ono w ustalonej
kostce $[-R,R]^k$ dla pewnego $R>0$. Dla $r\le1$
liczba kostek siatki przecinających $K$ nie przekracza
$(2R/r+2)^k\le C_{k,R}r^{-k}$. Łącząc powyższe
pokrycia, otrzymujemy pokrycie $G(K)$ zbiorami
o średnicach co najwyżej $C_k\eta r$ i koszcie
(co uwzględnia stałą normalizującą $c_k$)
nie większym niż $C_{k,N,M,R}\eta$.
Dla każdej zadanej skali $\delta>0$ możemy wybrać
$r$ tak małe, że $C_k\eta r\le\delta$ i wszystkie
wcześniejsze warunki są spełnione. Zatem
$\mathcal H^k_\delta(G(K))\le C_{k,N,M,R}\eta$.
Biorąc supremum po $\delta>0$, a następnie
przechodząc $\eta\downarrow0$, otrzymujemy
$\mathcal H^k(G(K))=0$.
\end{proof}

\begin{lemat}[Zbiory zerowe i ich włókna]
\label{lem:null-fibres-lipschitz-28}
Niech $n\ge m\ge1$, $F:\mathbb R^n\to\mathbb R^m$
będzie lokalnie Lipschitzowska, a
$Z\subseteq\mathbb R^n$ zbiorem o zerowej zewnętrznej
mierze Lebesgue'a. Wówczas
\[
 \mathcal H^{n-m}(Z\cap F^{-1}(\{y\}))=0
 \quad\text{dla }\lambda_m\text{-prawie wszystkich }y.
\]
\end{lemat}
\begin{proof}
Przez przeliczalne pokrycie dziedziny otwartymi
kostkami o domknięciach zawartych w większych otoczeniach,
na których $F$ jest Lipschitzowska, wystarczy rozpatrzyć
część $Z$ zawartą w jednej takiej kostce: dla każdej
części otrzymamy zbiór wyjątkowy o zewnętrznej mierze
zero, a ich przeliczalna suma nadal będzie zerowa.
Poza nią podaddytywność $\mathcal H^{n-m}$ da tezę
dla całego $Z$. Wybieramy dodatnią
stałą Lipschitza $L>0$ na większym otoczeniu (również
gdy $F$ jest na nim stałe).
Piszemy $d=n-m$. Ustalmy $\delta,\varepsilon>0$.
Z warunku $\lambda_n^*(Z)=0$ i kalibracji
$\mathcal H^n=\lambda_n$ (na zbiorach mierzalnych,
w szczególności zerowych) wynika $\mathcal H^n(Z)=0$.
Stąd dla dowolnej dodatniej skali $\tau$ istnieje
przeliczalne pokrycie $Z$ niepustymi zbiorami
$D_i$ o średnicach co najwyżej $\tau$ i sumie
$n$-tych potęg średnic dowolnie małej. Odrzucamy
zbiory nieprzecinające $Z$ i wybieramy $x_i\in D_i\cap Z$.
Jeśli $a_i=\operatorname{diam}D_i>0$, domknięta
kostka o środku $x_i$ i boku $2a_i$ pokrywa $D_i$,
a jej średnica wynosi $2\sqrt n\,a_i$.
Jeśli $a_i=0$, zbiór jest singletonem; pokrywamy
go kostką o dodatnim boku, wybranym tak małym,
aby jej średnica była mniejsza niż $2\sqrt n\,\tau$
i jej koszt nie przekraczał $\varepsilon 2^{-i-1}$.
Wybierając pierwotną sumę $n$-tych potęg średnic
mniejszą niż $\varepsilon/(2(2\sqrt n)^n)$,
otrzymujemy kostki $Q_i$ z
$\sum_i(\operatorname{diam}Q_i)^n<\varepsilon$.
Możemy wybrać $\tau$ tak małe, aby średnice były
mniejsze niż $\delta$ oraz niż odległość domknięcia
rozpatrywanej kostki od dopełnienia większego otoczenia,
na którym działa $L$ (jeżeli dopełnienie jest puste,
wystarczy pierwszy warunek). Wszystkie $Q_i$ leżą więc
w tym otoczeniu.
Niech $z_i$ będzie środkiem $Q_i$ i niech
$V_i=B(F(z_i),L\operatorname{diam}Q_i)$.
Jeśli $x\in Z\cap Q_i$ i $F(x)=y$, to
$|y-F(z_i)|\le L|x-z_i|\le L\operatorname{diam}Q_i/2
<L\operatorname{diam}Q_i$, a więc $y\in V_i$.
Dla ustalonego $y$ kostki o indeksach spełniających
$y\in V_i$ pokrywają całe włókno $Z\cap F^{-1}(\{y\})$.
Ich średnice są mniejsze niż $\delta$, dlatego
z definicji kosztu pokrycia otrzymujemy punktowo
\[
 \mathcal H^d_\delta(Z\cap F^{-1}(\{y\}))
 \le C_{n,m}\sum_i(\operatorname{diam}Q_i)^d\chi_{V_i}(y).
\]
Dla $d=0$ koszt niepustego zbioru w pokryciu
wynosi $1$, zgodnie z prawą stroną dla $C_{n,m}\ge1$.
Dla $d>0$ wystarcza $C_{n,m}\ge c_d$.
Prawa strona jest borelowska. Z twierdzenia
o zbieżności monotonicznej, ze wzoru
$\lambda_m(V_i)=\lambda_m(B(0,1))L^m
(\operatorname{diam}Q_i)^m$ i z $d+m=n$ otrzymujemy
\[
 \int_{\mathbb R^m}C_{n,m}\sum_i
 (\operatorname{diam}Q_i)^d\chi_{V_i}(y)
 \,\mathrm d\lambda_m(y)
 \le C_{n,m}L^m\varepsilon.
\]
Zatem dla dowolnego $t>0$ zbiór, na którym
$\mathcal H^d_\delta$ rozpatrywanego włókna
przekracza $t$, ma zewnętrzną miarę nie większą
niż $C_{n,m}L^m\varepsilon/t$. Ponieważ
$\varepsilon$ jest dowolne, miara ta jest zerowa.
Stosujemy to do $\delta=1/j$ i $t=1/l$ dla
$j,l\ge1$. Przeliczalna suma otrzymanych zbiorów
wyjątkowych ma zewnętrzną miarę zero. Poza nią
dla każdego $j$ zawartość $\mathcal H^d_{1/j}$
rozpatrywanego włókna nie przekracza $1/l$ dla
każdego $l$, więc jest zerowa. Przechodząc
$j\to\infty$, otrzymujemy $\mathcal H^d=0$
na prawie wszystkich włóknach.
\end{proof}

\begin{lemat}[Włókna zwartego zbioru krytycznego]
\label{lem:critical-coarea-null-28}
Niech $n\ge m\ge1$, $V\subseteq\mathbb R^n$
będzie otwarty, $G\in C^1(V;\mathbb R^m)$,
a $K\subseteq V$ zwarty. Jeżeli
$\operatorname{rank}DG(x)<m$ dla wszystkich $x\in K$,
to
\[
 \mathcal H^{n-m}(K\cap G^{-1}(\{y\}))=0
 \quad\text{dla }\lambda_m\text{-prawie wszystkich }y.
\]
\end{lemat}
\begin{proof}
Piszemy $d=n-m$. Dla ustalonej skali $\delta$ nie
zakładamy mierzalności funkcji
$y\mapsto\mathcal H^d_\delta(K\cap G^{-1}(\{y\}))$.
W definicji z Rozdziału~\ref{ch:hausdorff-26} dopuszcza
się dowolne zbiory o średnicy \emph{co najwyżej}
$\delta$; zastąpienie ich otwartymi zbiorami o średnicy
ściśle mniejszej niż $\delta$ może zmienić infimum.
Wystarczy oszacowanie zewnętrznej miary zbiorów
parametrów, jak w Lemacie~\ref{lem:null-fibres-lipschitz-28}.

Jeśli $K=\varnothing$, wszystkie włókna są puste.
Niech $K\ne\varnothing$ i $0<\eta\le1$.
Pokryjmy $K$ domkniętymi kostkami siatki
o boku $r$, dla których wybieramy $x_Q\in K\cap Q$.
Wybieramy zwarte otoczenie $K$ zawarte w $V$,
zawierające wszystkie punkty odległe od $K$
o co najwyżej pewną dodatnią liczbę.
Dla dostatecznie małego $r$ odcinki $[x_Q,x]$
dla $x\in K\cap Q$ leżą w tym otoczeniu,
a jednostajna ciągłość $DG$ zapewnia na nich
$\|DG(z)-DG(x_Q)\|\le\eta$.
Całkując różniczkę wzdłuż odcinka, dostajemy
\[
 |G(x)-G(x_Q)-DG(x_Q)(x-x_Q)|\le C_n\eta r
 \quad(x\in K\cap Q).
\]
Ponieważ ranga $DG(x_Q)$ jest mniejsza niż $m$,
zbiór $G(K\cap Q)$ zawiera się w
$C_n\eta r$-otoczeniu afinicznej podprzestrzeni
wymiaru $j=\operatorname{rank}DG(x_Q)\le m-1$,
w obrębie kuli o promieniu $C_n(M+1)r$, gdzie
$M$ ogranicza normę $DG$ na wybranym zwartym otoczeniu.
Po przesunięciu o $G(x_Q)$ i ortogonalnej zmianie
współrzędnych obraz jest zawarty w produkcie
$j$-wymiarowej kuli o promieniu $C_n(M+1)r$
i $(m-j)$-wymiarowej kuli o promieniu $C_n\eta r$.
Przy $j=0$ pierwszy czynnik jest jednopunktową
przestrzenią o zerowym wymiarze i mierze liczącej $1$.
Objętość produktu nie przekracza
$C_{n,m,M}r^j(\eta r)^{m-j}$, a
$\eta^{m-j}\le\eta$ dla $0<\eta\le1$.
Ponieważ $K\cap Q$ jest zwarty i $G$ ciągłe,
jego obraz jest zwarty, a zatem borelowski. Otrzymujemy
\[
 \lambda_m(G(K\cap Q))\le C_{n,m,M}\eta r^m.
\]
Dla każdego $y$ zbiór $K\cap G^{-1}(\{y\})$
można pokryć fragmentami leżącymi w tych
kostkach $Q$, dla których $y\in G(K\cap Q)$.
Każdy fragment ma średnicę nie większą niż
$\sqrt n\,r$, więc dla $\delta_r=2\sqrt n\,r$
\[
 \mathcal H^d_{\delta_r}(K\cap G^{-1}(\{y\}))
 \le C_{n,m}\sum_Q r^d\chi_{G(K\cap Q)}(y).
\]
W sumie występuje najwyżej $C_K r^{-n}$ kostek
dla $0<r\le1$, ponieważ $K$ leży w ustalonej
ograniczonej kostce. Oznaczmy prawą stronę
nierówności przez $b_r(y)$. Jest to nieujemna
funkcja borelowska. Z oszacowania objętości
obrazów i z równości $d+m=n$ wynika
\[
 \int_{\mathbb R^m}b_r\,\mathrm d\lambda_m
 \le C_{n,m,M,K}r^{-n}r^d\eta r^m
 =C_{n,m,M,K}\eta.
\]
Ustalmy teraz skalę $\delta>0$ i próg $t>0$.
Dla każdego $0<\eta\le1$ możemy dobrać $r$
tak małe, że powyższe oszacowania obowiązują
oraz $\delta_r\le\delta$.
Z monotoniczności zawartości względem skali mamy
$\mathcal H^d_\delta(K\cap G^{-1}(\{y\}))
\le\mathcal H^d_{\delta_r}(K\cap G^{-1}(\{y\}))
\le b_r(y)$.
W konsekwencji
\[
 \lambda_m^*\{y:\mathcal H^d_\delta
 (K\cap G^{-1}(\{y\}))>t\}
 \le\lambda_m(\{b_r>t\})
 \le C_{n,m,M,K}\eta/t.
\]
Stała po prawej nie zależy od $\eta,r,\delta,t$.
Przechodząc $\eta\downarrow0$, otrzymujemy
zewnętrzną miarę zero dla ustalonych $\delta,t$.
Bierzemy przeliczalną sumę tych zbiorów dla
$\delta=1/j$ oraz $t=1/l$, $j,l\ge1$.
Poza nią wszystkie zawartości
$\mathcal H^d_{1/j}(K\cap G^{-1}(\{y\}))$ są zerowe.
Ich granica dla $j\to\infty$ jest z definicji
miarą $\mathcal H^d$ włókna i również wynosi zero.
To daje tezę bez całkowania funkcji o nieustalonej
mierzalności.
\end{proof}


\section{Wzory area i coarea dla odwzorowań Lipschitzowskich}
W poniższych wzorach $DF$ oznacza różniczkę z twierdzenia
Rademachera, istniejącą prawie wszędzie. Jakobian $J_kF$
oraz jakobian normalny $\operatorname{Jac}_mF$ obliczamy
w punktach różniczkowalności, podstawiając $DF(x)$
do funkcji jakobianowych z
Definicji~\ref{def:k-jacobian-28} i
Definicji~\ref{def:normal-jacobian-28},
a na pozostałym zbiorze zerowym przypisujemy im wartość $0$.
Zbiór punktów różniczkowalności jest borelowski
na mocy dowodu Twierdzenia~\ref{thm:rademacher-28},
a współczynniki $DF$ są na nim granicami borelowskich
ilorazów różnicowych w kierunkach wektorów bazowych.
Funkcje jakobianowe macierzy są ciągłe, więc opisane
jakobiany są borelowskie. Wybór wartości na zbiorze
wyjątkowym nie zmienia całek. W całkach przyjmujemy
$0\cdot(+\infty)=0$. Sumę nieujemnych liczb po
dowolnym, także nieprzeliczalnym zbiorze indeksów,
rozumiemy jako supremum sum po jego skończonych
podzbiorach (suma pusta wynosi zero).
\begin{twierdzenie}[Wzory area i coarea dla odwzorowań Lipschitzowskich]
\label{thm:lipschitz-area-coarea-28}
Niech $1\le k\le N$, $E\subseteq\mathbb R^k$ będzie
mierzalny w sensie Lebesgue'a, a $F:\mathbb R^k\to\mathbb R^N$
będzie lokalnie Lipschitzowski.
Wówczas różniczka $DF(x)$ istnieje dla prawie
wszystkich $x$ (twierdzenie Rademachera)
i dla każdej funkcji $f:E\to[0,+\infty]$
mierzalnej w sensie Lebesgue'a zachodzi
\[
 \int_E f(x)J_kF(x)\,\mathrm d\lambda_k(x)
 =\int_{\mathbb R^N}\left(\sum_{x\in E\cap F^{-1}(\{y\})}
 f(x)\right)\mathrm d\mathcal H^k(y).
\]
Funkcja pod całką po $\mathbb R^N$ jest mierzalna
względem zupełnej miary Hausdorffa.
Ponadto jeżeli $n\ge m\ge1$ i
$F:\mathbb R^n\to\mathbb R^m$ jest lokalnie
Lipschitzowski, to dla każdej funkcji
$g:\mathbb R^n\to[0,+\infty]$ mierzalnej
w sensie Lebesgue'a jej ograniczenie do włókna
jest $\mathcal H^{n-m}$-mierzalne dla
$\lambda_m$-prawie wszystkich $y$ i zachodzi
\[
 \int_{\mathbb R^n}g(x)\operatorname{Jac}_mF(x)
 \,\mathrm d\lambda_n(x)
 =\int_{\mathbb R^m}\left(
 \int_{F^{-1}(\{y\})}g(x)
 \,\mathrm d\mathcal H^{n-m}(x)\right)
 \,\mathrm d\lambda_m(y).
\]
Wewnętrzna całka określa funkcję mierzalną względem
miary Lebesgue'a po przypisaniu dowolnej wartości
nieujemnej na ewentualnym zbiorze wyjątkowym miary zero.
W przypadku $n=m$ wewnętrzna całka po włóknie
jest sumą względem miary liczącej.
\end{twierdzenie}

\begin{lemat}[Wzór area dla dowolnego odwzorowania $C^1$]
\label{lem:c1-general-area-28}
Niech $1\le k\le N$, $V\subseteq\mathbb R^k$ będzie
otwarty, $G\in C^1(V;\mathbb R^N)$, a $A\subseteq V$
borelowski. Dla każdej nieujemnej funkcji borelowskiej
$f:A\to[0,+\infty]$ zachodzi
\[
 \int_Af(x)J_kG(x)\,\mathrm d\lambda_k(x)
 =\int_{\mathbb R^N}\left(
 \sum_{x\in A\cap G^{-1}(\{y\})}f(x)\right)
 \,\mathrm d\mathcal H^k(y).
\]
\end{lemat}
\begin{proof}
Zbiór $V_{\mathrm{reg}}=\{x\in V:\operatorname{rank}DG(x)=k\}$
jest otwarty, a $G|_{V_{\mathrm{reg}}}$ jest immersją.
Stosujemy Twierdzenie~\ref{thm:immersion-area-28}
do $G|_{V_{\mathrm{reg}}}$, zbioru borelowskiego
$A\cap V_{\mathrm{reg}}$ i ograniczenia $f$ do tego zbioru.
Jego dopełnienie w $V$ jest względnie domkniętym
zbiorem punktów krytycznych, który można wyczerpać
przeliczalną rodziną zwartych podzbiorów.
Istotnie, wystarczy przeciąć zbiór krytyczny ze
zwartymi zbiorami wyczerpującymi $V$, np. kulami
domkniętymi o promieniach $l$ ograniczonymi dodatkowo
warunkiem $\operatorname{dist}(x,\mathbb R^k\setminus V)
\ge1/l$, z pominięciem tego warunku gdy $V=\mathbb R^k$.
Każdy taki zbiór jest zwarty i spełnia hipotezy
Lematu~\ref{lem:critical-area-null-28}. Podaddytywność
pokazuje, że obraz całego zbioru krytycznego ma
$\mathcal H^k$-miarę zero.
Na części krytycznej $J_kG=0$, a suma po włóknach
może być różna od zera tylko na wspomnianym
zbiorze zerowym. Jest ona mierzalna względem zupełnej
miary Hausdorffa, ponieważ każdy podzbiór zbioru
zerowego jest mierzalny. Obie całki części krytycznej
wynoszą więc zero. Dodając równość dla części
regularnej, otrzymujemy żądany wzór.
\end{proof}

\begin{proof}[Dowód ogólnych wzorów area i coarea z
Twierdzenia~\ref{thm:lipschitz-area-coarea-28}]
\emph{Krok 1: wzór area.}
Z Lematu~\ref{lem:lusin-whitney-28} otrzymujemy
zwarte zbiory $K_j\subseteq\mathbb R^k$ oraz
odwzorowania $F_j\in C^1(\mathbb R^k;\mathbb R^N)$
takie, że $F_j=F$ i $DF_j=DF$ na $K_j$, a
$Z=\mathbb R^k\setminus\bigcup_jK_j$ jest
$\lambda_k$-zerowy. Połóżmy
$A_1=K_1$ i
$A_j=K_j\setminus\bigcup_{i<j}K_i$ dla $j\ge2$.
Są to parami rozłączne zbiory borelowskie,
których suma jest dopełnieniem $Z$.

Dla $E$ i $f$ borelowskich stosujemy
Lemat~\ref{lem:c1-general-area-28} do
$F_j$ na borelowskim zbiorze $A_j\cap E$
z ograniczeniem $f$ do tego zbioru. Ponieważ
$A_j\subseteq K_j$, dla każdego $x\in A_j$ mamy
$J_kF_j(x)=J_kF(x)$ oraz $F_j(x)=F(x)$.
Zatem dla każdego $y$ zachodzi
$A_j\cap E\cap F_j^{-1}(\{y\})
=A_j\cap E\cap F^{-1}(\{y\})$.
Podstawiając obie równości do wzoru z lematu, dostajemy
\[
 \int_{A_j\cap E}f(x)J_kF(x)\,\mathrm d\lambda_k(x)
 =\int_{\mathbb R^N}
 \sum_{x\in A_j\cap E\cap F^{-1}(\{y\})}f(x)
 \,\mathrm d\mathcal H^k(y).
\]
Po lewej stronie sumujemy całki po rozłącznych
zbiorach $A_j\cap E$. Po prawej każda funkcja
sumy po włóknie jest nieujemna i mierzalna
względem zupełnej miary Hausdorffa na mocy dowodu
lematu. Dla ustalonego $y$ połóżmy
$T_j=A_j\cap E\cap F^{-1}(\{y\})$.
Zbiory $T_j$ są rozłączne. Każdy skończony podzbiór
ich sumy przecina tylko skończenie wiele $T_j$,
więc jego suma wartości $f$ nie przekracza
$\sum_j\sum_{x\in T_j}f(x)$.
Biorąc supremum po skończonych podzbiorach, otrzymujemy
jedną nierówność. Odwrotnie, dla ustalonego $l$
i skończonych podzbiorów $D_j\subseteq T_j$,
$j=1,\ldots,l$, suma wartości $f$ na ich rozłącznej
sumie nie przekracza sumy po całym włóknie
z pominięciem $Z$. Biorąc kolejno supremum po
$D_1,\ldots,D_l$, a następnie po $l$, otrzymujemy
nierówność odwrotną. To dowodzi równości sum
także wtedy, gdy któraś ma wartość $+\infty$. Stosujemy twierdzenie
o zbieżności monotonicznej do sum częściowych,
które nie wymaga $\sigma$-skończoności $\mathcal H^k$
na $\mathbb R^N$.
Zbiór $Z$ nie wnosi nic do lewej strony.
Ponieważ $F$ jest lokalnie Lipschitzowskie,
wybieramy przeliczalne pokrycie $\mathbb R^k$
otwartymi kulami $U_i$, na których $F$ jest
Lipschitzowskie ze stałą $L_i>0$. Dla każdego $i$
kalibracja daje $\mathcal H^k(Z\cap U_i)=0$,
a Twierdzenie~\ref{thm:hausdorff-lipschitz-26} daje
$\mathcal H^k(F(Z\cap U_i))\le
L_i^k\mathcal H^k(Z\cap U_i)=0$.
Podaddytywność daje $\mathcal H^k(F(Z))=0$.
Zatem również prawa strona nie zmienia się
po dołączeniu punktów z $Z$.

Jeżeli $E$ i $f$ są jedynie mierzalne w sensie
Lebesgue'a, przedłużamy $f$ przez zero poza $E$
do funkcji $u:\mathbb R^k\to[0,+\infty]$.
Wybieramy borelowski reprezentant $u_0$ funkcji $u$
równy jej poza zbiorem $\lambda_k$-zerowym $D$.
Stosujemy dotychczasowy wzór do $E=\mathbb R^k$
i funkcji $u_0$. Całki po dziedzinie dla $u,u_0$
są równe. Jak wykazano wyżej, $\mathcal H^k(F(D))=0$.
Dla $y\notin F(D)$ wszystkie wartości $u$ i $u_0$
na włóknie są identyczne, a zerowe wartości poza $E$
nie wpływają na sumę. Zupełność miary Hausdorffa
zapewnia mierzalność rzeczywistej funkcji sumy
i równość całek po $\mathbb R^N$.

\emph{Krok 2: wzór coarea na częściach pełnej rangi.}
Przyjmijmy na razie, że $g$ jest borelowska.
Dla odwzorowania $F:\mathbb R^n\to\mathbb R^m$
stosujemy Lemat~\ref{lem:lusin-whitney-28},
otrzymując zwarte zbiory $K_j\subseteq\mathbb R^n$
i odwzorowania $F_j\in C^1(\mathbb R^n;\mathbb R^m)$
takie, że $F_j=F$ i $DF_j=DF$ na $K_j$, a
$\lambda_n(\mathbb R^n\setminus\bigcup_jK_j)=0$.
Kładziemy $A_j=K_j\setminus\bigcup_{i<j}K_i$,
więc są to parami rozłączne zbiory borelowskie.
Niech
$V_j=\{x:\operatorname{rank}DF_j(x)=m\}$.
Jest to zbiór otwarty, a $F_j|_{V_j}$ jest
submersją. Twierdzenie~\ref{thm:coarea-submersion-28}
stosujemy do $F_j|_{V_j}$ i nieujemnej borelowskiej
funkcji $(g\chi_{A_j})|_{V_j}$. Po lewej stronie
na $A_j\cap V_j$ jakobiany $F_j$ i $F$ są równe,
gdyż $DF_j=DF$ na $K_j$. Po prawej stronie
funkcja znika poza $A_j$, a na tym zbiorze
$F_j=F$, więc
$A_j\cap V_j\cap F_j^{-1}(\{y\})
=A_j\cap V_j\cap F^{-1}(\{y\})$.
Wzór dla submersji daje zatem
\begin{align*}
 &\int_{A_j\cap V_j}g(x)\operatorname{Jac}_mF(x)
 \,\mathrm d\lambda_n(x)\\
 &\quad=\int_{\mathbb R^m}\left(
 \int_{A_j\cap V_j\cap F^{-1}(\{y\})}
 g(x)\,\mathrm d\mathcal H^{n-m}(x)\right)
 \,\mathrm d\lambda_m(y).
\end{align*}

\emph{Krok 3: części krytyczne i zbiór wyjątkowy.}
Zbiór
$C_j=K_j\cap\{x:\operatorname{rank}DF_j(x)<m\}$
jest zwarty, ponieważ warunek rangi mniejszej niż $m$
jest domknięty (wszystkie minory rzędu $m$ znikają).
Na $C_j\subseteq K_j$ mamy $DF_j=DF$, więc
$\operatorname{Jac}_mF=0$. Lemat~\ref{lem:critical-coarea-null-28}
stosujemy do odwzorowania $F_j$ i zbioru $C_j$:
wszystkie punkty $C_j$ są krytyczne dla $F_j$.
Ponieważ $F_j=F$ na $C_j$, otrzymujemy
$\mathcal H^{n-m}(C_j\cap F^{-1}(\{y\}))=0$
dla prawie wszystkich $y$. Na pozostałym zbiorze
$Z=\mathbb R^n\setminus\bigcup_jK_j$ jakobian
również nie wnosi nic do całki, bo
$\lambda_n(Z)=0$. Z kolei
Lemat~\ref{lem:null-fibres-lipschitz-28} dowodzi,
że $Z$ jest $\mathcal H^{n-m}$-zerowy na
prawie wszystkich włóknach. Bierzemy przeliczalną sumę zbiorów wyjątkowych
dla wszystkich $C_j$ i dla $Z$. Poza nią brakujące
części włókna mają zerową $\mathcal H^{n-m}$-miarę.
Dla borelowskiego $g$ ograniczenia do włókien są
mierzalne, ponieważ $F$ jest ciągłe, włókna domknięte,
a miara Hausdorffa mierzy zbiory borelowskie.
Dla każdego $y$ spoza zbioru wyjątkowego
całka po włóknie jest sumą całek po rozłącznych
$A_j\cap V_j\cap F^{-1}(\{y\})$.
Każda funkcja tej lokalnej całki jest mierzalna
względem $\lambda_m$ na mocy dowodu wzoru
dla submersji (Tonelli i rozłączny podział).
Suma jest mierzalna, a rzeczywista całka po włóknie
różni się od niej co najwyżej na zbiorze zerowym;
stosujemy zupełność miary Lebesgue'a.
Po obu stronach sumujemy równości z kroku 2
przez twierdzenie o zbieżności monotonicznej.
Pominięte części dają całki zerowe, co dowodzi wzoru.

W przypadku funkcji $g$ jedynie mierzalnej
w sensie Lebesgue'a zastępujemy ją wersją
borelowską poza zbiorem zerowym. Niech $D$ będzie zbiorem zerowym, poza którym
$g$ i jego borelowski reprezentant $g_0$ są równe.
Lemat~\ref{lem:null-fibres-lipschitz-28} daje
$\mathcal H^{n-m}(D\cap F^{-1}(\{y\}))=0$
dla prawie wszystkich $y$. Zupełność miary Hausdorffa
zapewnia dla tych $y$ mierzalność ograniczenia $g$
do włókna i równość jego całki z całką $g_0$.
Na pozostałych włóknach całce możemy przypisać
dowolną nieujemną wartość; zupełność miary Lebesgue'a
zapewnia mierzalność tak otrzymanej funkcji i równość
jej całki z całką odpowiadającą $g_0$.
Nie trzeba w żadnym kroku zakładać, że odwzorowanie
$F$ jest globalnie różnowartościowe lub że ma
stałą rangę.
\end{proof}

\begin{uwaga}[Zakres uzyskanego dowodu]
W porównaniu z wersjami gładkimi nowym elementem
jest jednoczesna kontrola zbiorów krytycznych i ich
włókien. Lemat~\ref{lem:whitney-c1-28} wyprowadza
potrzebne przedłużenie $C^1$, a lematy
\ref{lem:critical-area-null-28}--\ref{lem:critical-coarea-null-28}
wyprowadzają oszacowania bez powoływania się na
same wzory area/coarea. Dzięki temu dowód nie jest
błędnym kołem. Nadal poza zakresem pozostają
analogiczne wzory na ogólnych przestrzeniach
metrycznych oraz dla odwzorowań, które nie są
lokalnie Lipschitzowskie.
\end{uwaga}

\section{Postać na rozmaitościach}
\begin{twierdzenie}[Wzór koobszaru na rozmaitościach riemannowskich]
Niech $(M^n,g)$ i $(N^m,h)$ będą gładkimi rozmaitościami
riemannowskimi bez brzegu, Hausdorffa, spełniającymi drugi
aksjomat przeliczalności, wymiarów $n\ge m\ge1$.
Niech $F:M\to N$ będzie submersją klasy $C^1$, tj.
$\mathrm dF_p:T_pM\to T_{F(p)}N$ jest surjekcją dla każdego
$p\in M$. Przez $\mathrm dF_p^*$ oznaczamy sprzężenie
różniczki względem iloczynów skalarnych $g_p,h_{F(p)}$.
Jakobian normalny definiujemy jako
$\operatorname{Jac}_{g,h}F(p)=
\sqrt{\det(\mathrm dF_p\mathrm dF_p^*)}$,
gdzie wyznacznik liczymy na $T_{F(p)}N$.
Dla każdej funkcji borelowskiej $\varphi:M\to[0,+\infty]$
funkcja wewnętrznej całki w poniższym wzorze jest
borelowska i zachodzi
\begin{equation}\label{eq:coarea-manifolds-28}
 \int_M\varphi(p)\operatorname{Jac}_{g,h}F(p)
 \,\mathrm d\mu_g(p)
 =\int_N\left(\int_{F^{-1}(\{q\})}\varphi(p)
 \,\mathrm d\mu_{g|_{F^{-1}(\{q\})}}(p)\right)
 \,\mathrm d\mu_h(q).
\end{equation}
Na niepustym włóknie, będącym osadzoną podrozmaitością
klasy $C^1$, miarę indukowaną definiujemy za pomocą
ciągłej dodatnio określonej macierzy metryki w lokalnych
parametryzacjach klasy $C^1$. Dla $n=m$ jest to miara
licząca, a całka po pustym włóknie wynosi zero.
\end{twierdzenie}
\begin{proof}
\emph{Krok 1: współrzędne dostosowane do submersji.}
Ustalmy $p_0\in M$ i mapy $u$ wokół $p_0$ oraz $v$ wokół
$F(p_0)$. Ograniczamy dziedzinę $u$, aby jej obraz przez $F$
mieścił się w dziedzinie $v$. Odwzorowanie
$T=v\circ F\circ u^{-1}$ jest klasy $C^1$, a jego różniczka
w $u(p_0)$ ma rangę $m$. Po permutacji współrzędnych
źródła piszemy $a=(a',a'')\in\mathbb R^m\times\mathbb R^{n-m}$
tak, aby minor $D_{a'}T$ był odwracalny w $u(p_0)$.
Różniczka odwzorowania $a\mapsto(T(a),a'')$
ma wtedy postać blokowo trójkątną z odwracalnymi
blokami diagonalnymi. Z twierdzenia o odwzorowaniu
odwrotnym otrzymujemy jego lokalną odwrotność
$\Psi:W\to u(V)$ klasy $C^1$, gdzie $W$ jest otwarty
w $\mathbb R^m\times\mathbb R^{n-m}$, a $V$ jest
otwartym otoczeniem $p_0$ w $M$.
Kładziemy $P=u^{-1}\circ\Psi:W\to V$.
Dla $(y,z)\in W$ mamy $v(F(P(y,z)))=y$.
W szczególności dla $q$ z dziedziny mapy $v$,
przy $y=v(q)$ i $W_y=\{z:(y,z)\in W\}$,
odwzorowanie $z\mapsto P(y,z)$ parametryzuje
$V\cap F^{-1}(\{q\})$. Gdy $W_y$ jest puste,
ten fragment włókna jest pusty.

\emph{Krok 2: tożsamość macierzy metryk.}
W punkcie $(y,z)\in W$ niech $B$ oznacza macierz
metryki $g$ w parametryzacji $P$, czyli
$B=D\Psi^{\mathsf T}G_u(\Psi)D\Psi$, z notacją
macierzy $G_u$ z \eqref{eq:g-density-27}.
Niech $H=H_v(y)$ będzie macierzą metryki $h$
w mapie $v$. Piszemy
$B=\begin{psmallmatrix}B_{11}&B_{12}\\B_{21}&C\end{psmallmatrix}$,
gdzie $C$ jest blokiem rozmiaru $d\times d$, $d=n-m$.
Dla $d>0$ jest on dodatnio określony i jest właśnie
macierzą metryki indukowanej w parametryzacji włókna.
W tych współrzędnych różniczka $F$ ma macierz
$L=(I_m\ \ 0)$. Macierz jej sprzężenia wynosi
$B^{-1}L^{\mathsf T}H$: dla dowolnych wektorów
$a\in\mathbb R^n,b\in\mathbb R^m$ zachodzi
$(La)^{\mathsf T}Hb=a^{\mathsf T}B(B^{-1}L^{\mathsf T}Hb)$.
Zatem kwadrat jakobianu normalnego wynosi
$\det(LB^{-1}L^{\mathsf T})\det H$.

Dla $d>0$ kładziemy $S=B_{11}-B_{12}C^{-1}B_{21}$.
Mnożąc blokowo, otrzymujemy
\[
 \begin{pmatrix}I_m&-B_{12}C^{-1}\\0&I_d\end{pmatrix}
 B
 \begin{pmatrix}I_m&0\\-C^{-1}B_{21}&I_d\end{pmatrix}
 =\begin{pmatrix}S&0\\0&C\end{pmatrix}.
\]
Macierze mnożące mają wyznacznik $1$, skąd
$\det B=\det S\det C$. Dodatnia określoność $B$
i powyższa kongruencja dają dodatnią określoność $S$.
Jeśli $R$ oznacza prawą macierz mnożącą,
to lewa macierz mnożąca jest równa $R^{\mathsf T}$.
Zatem $R^{\mathsf T}BR=\operatorname{diag}(S,C)$
i po odwróceniu oraz pomnożeniu przez $R$ i
$R^{\mathsf T}$ dostajemy
\[
 B^{-1}=R\operatorname{diag}(S^{-1},C^{-1})R^{\mathsf T}.
\]
Pierwszy wiersz blokowy $R$ jest równy $(I_m\ \ 0)$,
więc lewy górny blok $B^{-1}$ jest równy $S^{-1}$. Ponieważ $L$ wybiera
ten blok, mamy $\det(LB^{-1}L^{\mathsf T})=
1/\det S=\det C/\det B$.
Dla $d=0$ mamy $L=I_m$, a przy konwencji $\det C=1$
ta ostatnia równość również obowiązuje. Stąd w obu przypadkach
\[
 \operatorname{Jac}_{g,h}F(P(y,z))\sqrt{\det B}
 =\sqrt{\det H}\sqrt{\det C}.
\]
Wszystkie macierze metryk są ciągłe i dodatnio określone;
w szczególności jakobian normalny jest dodatni i ciągły.

\emph{Krok 3: całkowanie lokalne.}
Niech $\varphi$ znika poza $V$. Współrzędnościowy wzór
na $\mu_g$ i zamiana zmiennych przez dyfeomorfizm
$\Psi$ klasy $C^1$ dają gęstość
$\sqrt{\det G_u(\Psi)}|\det D\Psi|=\sqrt{\det B}$
w parametryzacji $P$. Analogicznie wzór na miarę
indukowaną we włóknie daje gęstość $\sqrt{\det C}$.
Ta definicja jest niezależna od parametryzacji:
przy zmianie parametrów przez dyfeomorfizm klasy $C^1$
macierz metryki zmienia się z $C$ na $D^{\mathsf T}CD$,
gdzie $D$ jest macierzą różniczki zmiany parametrów
(z $C$ obliczanym w odpowiednim punkcie).
Pierwiastek z wyznacznika mnoży się przez $|\det D|$,
a euklidesowa zamiana zmiennych usuwa tę zmianę.
Dla $d=0$ nie występuje zmiana gęstości: każdy punkt
ma miarę $1$.

Dla $d>0$ definiujemy na $W$ nieujemną funkcję borelowską
\[
 b(y,z)=\varphi(P(y,z))\sqrt{\det C(y,z)}.
\]
Przedłużamy $b$ przez zero na całe $\mathbb R^m\times
\mathbb R^d$. Korzystając z tożsamości z kroku 2
i twierdzenia Tonellego dla $b$ (z czynnikiem
$\sqrt{\det H_v(y)}$ w całce po $y$), otrzymujemy
\begin{align*}
 &\int_V\varphi(p)\operatorname{Jac}_{g,h}F(p)\,\mathrm d\mu_g(p)\\
 &\quad=\int_W\varphi(P(y,z))\sqrt{\det H_v(y)}
      \sqrt{\det C(y,z)}\,\mathrm d\lambda_n(y,z)\\
 &\quad=\int_{v(\operatorname{dom} v)}
       \sqrt{\det H_v(y)}
       \left(\int_{W_y}\varphi(P(y,z))
              \sqrt{\det C(y,z)}\,\mathrm d\lambda_d(z)\right)
       \mathrm d\lambda_m(y).
\end{align*}
Wewnętrzna całka jest borelowską funkcją $y$, gdyż
jest równa $\int_{\mathbb R^d}b(y,z)\,\mathrm d\lambda_d(z)$,
a $b$ jest borelowska i nieujemna oraz $\lambda_d$
jest $\sigma$-skończoną miarą borelowską.
Współrzędnościowy wzór na $\mu_h$ i definicja miary
indukowanej identyfikują ostatnią całkę z prawą stroną
\eqref{eq:coarea-manifolds-28} dla funkcji znikającej
poza $V$. Gdy $d=0$, $P$ jest lokalną odwrotnością
$F$ zapisaną w mapie $v$, a każdy fragment włókna w $V$
jest pusty albo jednopunktowy. Ta sama tożsamość
wyznaczników i zamiana zmiennych dają równość,
w której całka wewnętrzna jest wartością funkcji
w jedynym punkcie fragmentu albo zerem.

\emph{Krok 4: przejście do całego $M$.}
Z otoczeń $V$ skonstruowanych w kroku 1 wybieramy
przeliczalne podpokrycie $(V_j)_{j\ge1}$, korzystając
z drugiego aksjomatu przeliczalności. Zbiory
$B_j=V_j\setminus\bigcup_{i<j}V_i$ są borelowskie,
parami rozłączne i pokrywają $M$.
Stosujemy lokalny wzór do $\varphi\chi_{B_j}$.
Dla każdego $q\in N$ zbiory $B_j\cap F^{-1}(\{q\})$
tworzą rozłączny borelowski podział włókna, więc
całka po włóknie jest sumą lokalnych całek.
Każda funkcja lokalnej całki jest borelowska
(równa zero poza odpowiednią dziedziną mapy celu),
a ich suma również jest borelowska.
Twierdzenie o zbieżności monotonicznej pozwala
zsumować lokalne równości po obu stronach.
Otrzymujemy \eqref{eq:coarea-manifolds-28}.
Nie wykorzystaliśmy orientowalności.
\end{proof}

\begin{wniosek}[Miara riemannowska i Hausdorffa]
\label{cor:riemann-hausdorff-28}
Niech $(M^n,g)$ będzie gładką rozmaitością riemannowską
bez brzegu, Hausdorffa, spełniającą drugi aksjomat
przeliczalności, wymiaru $n\ge1$. Na każdej składowej
spójności $M$ niech $d_g$ oznacza odległość riemannowską,
a $\mathcal H^n_{d_g}$ miarę Hausdorffa skonstruowaną
z tej odległości przy kalibracji z
Uwagi~\ref{rem:hausdorff-convention-26}.
Na zbiorach borelowskich zachodzi
$\mu_g=\mathcal H^n_{d_g}$ (dla rozmaitości niespójnej
miary sumujemy po składowych spójności).
Jeżeli $S^k\subseteq M$ jest osadzoną gładką
podrozmaitością bez brzegu, $0\le k\le n$, to
$\mu_{g|_S}=\mathcal H^k_{d_g|_{S\times S}}$
na zbiorach borelowskich $S$, również z sumowaniem
po składowych spójności $M$. Dla $k=0$ obie
miary na $S$ rozumiemy jako miary liczące.
\end{wniosek}
\begin{proof}
Wystarczy pracować na spójnej składowej $M$.
Składowe są otwarte, ponieważ rozmaitość jest
lokalnie spójna drogowo; jest ich co najwyżej
przeliczalnie wiele, gdyż każda zawiera inny element
przeliczalnej bazy. Odległość $d_g$ na każdej
składowej jest infimum długości odcinkami gładkich
krzywych łączących dane punkty.
Najpierw rozpatrujemy $1\le k\le n$, przy czym
przypadek $S=M$ odpowiada $k=n$.

\emph{Krok 1: lokalne porównanie odległości.}
Ustalmy $p\in S$ i $0<\varepsilon<1$. Wybieramy gładką
mapę $u$ otoczenia $p$ w $M$ taką, że $u(p)=0$
oraz $S$ w tej mapie odpowiada przecięciu jej obrazu
z podprzestrzenią $E=\mathbb R^k\times\{0\}$.
Gdy $S=M$, bierzemy $E=\mathbb R^n$.
Niech $G(a)=G_u(a)$ będzie macierzą metryki,
$G_0=G(0)$ oraz $|v|_0=\sqrt{v^{\mathsf T}G_0v}$.
Z ciągłości macierzy i dodatniej określoności $G_0$
wybieramy $R>0$, tak aby domknięta kula
$\{a:|a|_0\le R\}$ była zawarta w obrazie mapy
i dla $|a|_0\le R$, $v\in\mathbb R^n$ zachodziło
\[
 (1-\varepsilon)^2|v|_0^2
 \le v^{\mathsf T}G(a)v
 \le(1+\varepsilon)^2|v|_0^2.
\]
Istotnie, na zwartej sferze $|v|_0=1$ różnica
$v^{\mathsf T}(G(a)-G_0)v$ dąży jednostajnie do zera
dla $a\to0$; dla dowolnych $v$ stosujemy jednorodność.
Niech $0<r<R/3$ i $U_r=u^{-1}(\{a:|a|_0<r\})$.

Dla $x,y\in U_r$ odcinek między $u(x)$ i $u(y)$
pozostaje w kuli o promieniu $r$, ponieważ jest ona
wypukła. Jego przeciwobraz przez $u$ jest krzywą
łączącą $x,y$ o długości nie większej niż
$(1+\varepsilon)|u(x)-u(y)|_0$.
Z definicji infimum wynika takie samo górne
oszacowanie $d_g(x,y)$.

Aby otrzymać dolne oszacowanie, rozważamy dowolną
odcinkami gładką krzywą $\gamma$ łączącą $x,y$.
Jeśli pozostaje w przeciwobrazie kuli o promieniu
$R$, to jej długość jest co najmniej
$(1-\varepsilon)$ razy długość krzywej $u\circ\gamma$
w normie $|\cdot|_0$, a ta długość nie jest mniejsza
niż $|u(x)-u(y)|_0$ (nierówność trójkąta dla
całki z prędkości). Jeśli krzywa opuszcza tę kulę,
jej początkowy odcinek do pierwszego punktu na
sferze o promieniu $R$ leży w dziedzinie mapy:
domknięta kula jest w niej zawarta i jej przeciwobraz
jest zwarty. Długość tego odcinka jest co najmniej
$(1-\varepsilon)(R-r)$, gdyż w normie $|\cdot|_0$
odległość punktu początkowego od sfery wynosi
co najmniej $R-r$. Ponieważ $R-r>2r\ge|u(x)-u(y)|_0$,
również taka krzywa spełnia potrzebne dolne
oszacowanie. Biorąc infimum po wszystkich krzywych,
otrzymujemy
\[
 (1-\varepsilon)|u(x)-u(y)|_0
 \le d_g(x,y)
 \le(1+\varepsilon)|u(x)-u(y)|_0
 \quad(x,y\in U_r).
\]
Nie zakładamy istnienia krzywej minimalizującej.
Dla $x,y\in S\cap U_r$ górne oszacowanie pochodzi
nawet od krzywej w $S$, bo odcinek współrzędnych
leży w $E$. Dolne oszacowanie obowiązuje dla wszystkich
krzywych w $M$, a zatem dotyczy dokładnie odległości
zewnętrznej $d_g$ ograniczonej do $S$.

\emph{Krok 2: porównanie obu miar na małym otoczeniu.}
Przez $C(a)$ oznaczamy macierz ograniczenia $G(a)$
do $E$ w standardowej bazie $E$ i kładziemy $C_0=C(0)$.
Identyfikujemy $E$ z $\mathbb R^k$. Ograniczając
nierówność dla $G(a)$ do wektorów w $E$, otrzymujemy
oszacowanie wszystkich wartości własnych symetrycznej
macierzy $C_0^{-1/2}C(a)C_0^{-1/2}$ pomiędzy
$(1-\varepsilon)^2$ i $(1+\varepsilon)^2$.
Iloczyn tych wartości jest jej wyznacznikiem,
równym $\det C(a)/\det C_0$. Zatem
\[
 (1-\varepsilon)^k\sqrt{\det C_0}
 \le\sqrt{\det C(a)}
 \le(1+\varepsilon)^k\sqrt{\det C_0}.
\]
Wzór na gęstość miary indukowanej daje dla
borelowskiego $A\subseteq S\cap U_r$
\[
 (1-\varepsilon)^k D_A\le\mu_{g|_S}(A)
 \le(1+\varepsilon)^k D_A,\qquad
 D_A=\sqrt{\det C_0}\,\lambda_k(u(A)),
\]
gdzie w ostatniej mierze identyfikujemy $E$ z $\mathbb R^k$.
Dla $S=M$ jest to ten sam wzór z $C=G$.

Rozważmy na $E$ odległość $d_0(a,b)=|a-b|_0$.
Odwzorowanie $a\mapsto C_0^{1/2}a$ w bazie $E$
jest izometrią na zwykłą przestrzeń euklidesową.
Z definicji miary Hausdorffa wynika jej niezmienniczość
względem izometrii (pokrycia i średnice odpowiadają
sobie w obie strony). Kalibracja i liniowa zamiana
zmiennych dają więc $\mathcal H^k_{d_0}(u(A))=D_A$.
Miara Hausdorffa zbioru zawartego w podprzestrzeni
nie zależy od tego, czy pokrycia wybieramy w tej
podprzestrzeni, czy w przestrzeni otaczającej:
w drugim przypadku przecinamy każdy zbiór pokrycia
z podprzestrzenią, nie zwiększając średnic ani kosztów.
Zatem stosujemy Twierdzenie~\ref{thm:hausdorff-lipschitz-26}
do ograniczenia $u$ i jego odwrotności pomiędzy
$S\cap U_r$ z odległością $d_g$ oraz jego obrazem
z odległością $d_0$. Krok 1 daje odpowiednie stałe
Lipschitza $(1-\varepsilon)^{-1}$ i $1+\varepsilon$.
Otrzymujemy
\[
 (1-\varepsilon)^kD_A
 \le\mathcal H^k_{d_g|_{S\times S}}(A)
 \le(1+\varepsilon)^kD_A.
\]
Łącząc oba porównania (także dla $D_A=0$ lub
$D_A=+\infty$), dostajemy
\[
 \left(\frac{1-\varepsilon}{1+\varepsilon}\right)^k
 \mu_{g|_S}(A)
 \le\mathcal H^k_{d_g|_{S\times S}}(A)
 \le
 \left(\frac{1+\varepsilon}{1-\varepsilon}\right)^k
 \mu_{g|_S}(A).
\]

\emph{Krok 3: przejście do całej podrozmaitości.}
Dla ustalonego $\varepsilon$ otoczenia $S\cap U_r$
z kroku 1 pokrywają $S$. Podrozmaitość $S$ z topologią
podprzestrzeni spełnia drugi aksjomat przeliczalności,
więc wybieramy przeliczalne podpokrycie $(O_j)_{j\ge1}$.
Zbiory $O_j\setminus\bigcup_{i<j}O_i$ tworzą rozłączny
borelowski podział $S$. Obie miary mierzą zbiory
borelowskie w tej topologii: dla Hausdorffa wynika
to z konstrukcji miary na przestrzeni metrycznej
i faktu, że odległość $d_g$ indukuje topologię $M$
(co potwierdzają także lokalne porównania z kroku 1).
Dla każdego borelowskiego $A\subseteq S$ rozkładamy
$A$ według powyższego podziału i sumujemy nierówności
z kroku 2. Otrzymujemy te same globalne nierówności,
ze stałymi zależnymi jedynie od $k,\varepsilon$.
Jeśli $\mu_{g|_S}(A)<\infty$, przechodząc
$\varepsilon\downarrow0$, otrzymujemy równość obu miar.
Jeśli $\mu_{g|_S}(A)=+\infty$, już dolna nierówność
dla dowolnego $\varepsilon\in(0,1)$ daje nieskończoność
miary Hausdorffa. To kończy dowód dla $k\ge1$.
Dla $k=0$ osadzona podrozmaitość jest dyskretna
i przeliczalna. Miara indukowana nadaje każdemu
punktowi masę $1$, a $\mathcal H^0$ jest miarą
liczącą z definicji i wyniku z Rozdziału~\ref{ch:hausdorff-26}.
Obie miary są więc równe również w tym przypadku.
\end{proof}

\begin{uwaga}[Co przechodzi do geometrycznej teorii miary?]
Następny rozdział zastosuje miary Hausdorffa,
formuły area i coarea oraz twierdzenia o
różniczkowaniu miar do zbiorów prostowalnych
i zbiorów o skończonym obwodzie. Te obiekty
nie muszą być gładkimi podrozmaitościami,
a zatem stanowią właściwe rozszerzenie
geometrii różniczkowej.
\end{uwaga}

\begin{uwaga}[Dżety przybliżone a wzory miarowo-geometryczne]
Pochodna $DF(x)$ z twierdzenia Rademachera jest częścią
pierwszego dżetu odwzorowania $F$ w punktach różniczkowalności.
Przy osłabieniu pojęcia granicy do granicy przybliżonej otrzymujemy
\emph{dżety przybliżone}, określone miarowo prawie wszędzie.
W rozdziale~\ref{ch:dzety} rozwijamy geometryczną teorię wiązek
dżetów, omawiamy twierdzenie Whitneya o przedłużaniu danych
dżetowych oraz jego miarową wersję Whitneya--Łuzina.
Jest to uzupełnienie obecnego rozdziału, a nie ponowny dowód
wzorów area i coarea.
\end{uwaga}
